\documentclass[10pt]{article}
\usepackage[letterpaper,margin=1in]{geometry}
\usepackage[T1]{fontenc}
\usepackage{lmodern,amsmath,amssymb,amsthm,mathtools,enumitem,booktabs,graphicx}
\usepackage{microtype}
\usepackage[colorlinks=true,linkcolor=blue,citecolor=blue,urlcolor=blue]{hyperref}
\newtheorem{theorem}{Theorem}[section]
\newtheorem{proposition}[theorem]{Proposition}
\newtheorem{lemma}[theorem]{Lemma}
\newtheorem{corollary}[theorem]{Corollary}
\theoremstyle{remark}\newtheorem{remark}[theorem]{Remark}
\DeclareMathOperator{\arctanh}{arctanh}
\title{OOD Reversals in Two-Layer ReLU Networks:\\
Geometry, Timing, and Provable Non-Monotonicity\\Appendix}
\author{}\date{}
\begin{document}\maketitle
\noindent
This appendix separates the empirical geometry and timing protocol
(Appendix~\ref{app:empirical}), the pointwise strong-block normal form that
defines the empirical timing quantity $\mathcal G$
(Appendix~\ref{app:normal}), a computer-assisted canonical continuation and
timing theorem at the experimental parameters
(Appendix~\ref{app:canonical-cert}), additional general OOD results
(Appendix~\ref{app:general-theory}), and a finite-sample random-initialization
theorem proving OOD non-monotonicity in an explicit noisy regime
(Appendix~\ref{app:finite}). Appendix~\ref{app:validation} reports numerical
consistency checks at that certified parameter point.
Appendix~\ref{app:noiseless} gives an exact initialized trajectory on which
the normal-form interface is non-vacuous and its approximation error is
uniformly controlled.

\paragraph{Dependency map.}
The empirical claims have two distinct roles. Appendix~\ref{app:empirical}
first establishes where the observed reversals live and then tests the
zero crossing of $\mathcal G$ as a timing statistic. Appendix~\ref{app:normal}
defines that statistic and states the reduction error required for its sign
to control the full OOD derivative. Appendix~\ref{app:canonical-cert} then
uses a separate computer-assisted argument on the exact canonical dataset:
conditional on entry into a $1.2\cdot10^{-5}$ neighborhood of the learned
$t=3.4$ state, it certifies the full-network reversal on an explicit sector
and proves that the repository statistic $Q_G=\eta^2\mathcal G$ changes sign
strictly before every full OOD-error minimizer on the certified window.
Exact initialization-to-entry reachability remains open. Appendix~\ref{app:noiseless}
proves an exact initialized model in which the same normal form holds uniformly
along the true trajectory. By contrast, Appendix~\ref{app:finite} is
deliberately independent of the normal form: it proves initial OOD
improvement followed by later degradation directly from random
initialization in a different explicit noisy regime. Appendix~\ref{app:general-theory}
collects supporting structural certificates that are useful beyond either
headline result.
\setcounter{tocdepth}{1}
\tableofcontents
\clearpage\appendix
\section{Setup and common technical preliminaries}\label{app:setup}
\subsection{Model and notation}
For an even sample count $N$, the empirical two-concept identity-mapping
problem has $N/2$ independent samples per cluster,
$x_k^{(p)}=\mu_pe_p+\xi_k^{(p)}$, with target $x_k^{(p)}$.
Here $e_1,e_2$ are orthonormal. In the planar isotropic experiments and
Theorem~\ref{sbsharp:theorem}, $\xi_k^{(p)}\sim\mathcal N(0,\sigma^2I_2)$.
Write $\mathbb E_n$ for the average over all $N$ samples. The network and loss are
\[
 f_t(x)=\sum_{i=1}^h w_i(t)(u_i(t)^\top x)_+,
 \qquad \mathcal L(t)=\tfrac12\mathbb E_n\|f_t(x)-x\|^2.
\]
Rows of $U,W\in\mathbb R^{h\times d}$ are $u_i^\top,w_i^\top$.
Initialization is $u_i(0)=w_i(0)=\varepsilon\widehat u_i$, with independent
$\widehat u_i\sim\mathrm{Unif}(\mathbb S^{d-1})$, independent of the data.
Neuronwise balance, proved below, defines the mass
\[
 S(t)=\|U(t)\|_F^2=\|W(t)\|_F^2,\qquad S(0)=h\varepsilon^2=:s.
\]
The full squared parameter norm is $2S$.
Dots denote physical-time derivatives; primes in Appendix~\ref{app:finite}
denote derivatives in $\tau=t/2$. For a unit probe $\xi$,
$E(\xi,t)=\tfrac12\|f_t(\xi)-\xi\|^2$.
All time-derivative inequalities are almost-everywhere statements.

\subsection{Nonsmooth flow and balance}
\begin{lemma}[Clarke flow and compatible selectors]\label{lem:gf-regularity}
For any finite dataset, the loss is locally Lipschitz and semialgebraic.
From every finite initial state there exists a global locally absolutely
continuous solution of $-\dot\theta\in\partial^\circ\mathcal L(\theta)$,
where $\theta=(U,W)$ and $\partial^\circ$ is the Clarke subdifferential.
Every such solution admits measurable selectors $\alpha_i(t;x)\in[0,1]$,
equal to $1$ or $0$ at positive or negative training preactivations, such that
\begin{equation}\label{eq:flow}
 \dot w_i=\mathbb E_n[r_x(u_i^\top x)_+],\qquad
 \dot u_i=\mathbb E_n[\langle w_i,r_x\rangle\alpha_i(x)x],
 \qquad r_x=x-f_t(x).
\end{equation}
They satisfy $\alpha_i(x)u_i^\top x=(u_i^\top x)_+$.
For any fixed probe, the network output and its squared error are locally
absolutely continuous. The loss obeys
\begin{equation}\label{eq:dissipation}
 \dot{\mathcal L}=-\|\dot U\|_F^2-\|\dot W\|_F^2\le0.
\end{equation}
\end{lemma}
\begin{proof}
The finitely many signs of $u_i^\top x$ partition parameter space into
polyhedral cells. On each cell $f$ is polynomial and $\mathcal L$ is a
polynomial of degree at most four. Continuity across their boundaries
proves local Lipschitzness and semialgebraicity. The Clarke subdifferential
is nonempty, compact, convex, locally bounded and upper semicontinuous.
The local existence theorem for such differential inclusions (equivalently,
Euler polygonal approximation followed by compactness and the closed convex
graph limit) supplies a locally absolutely continuous trajectory.

At a differentiability point, ordinary differentiation gives
\eqref{eq:flow} with binary gates. At a boundary, every limiting gradient
has such a representation; a convex combination of limiting gradients
replaces each boundary gate by a number in $[0,1]$. The other coefficients
are continuous. Thus every Clarke element has a compatible representation.
Along a measurable velocity, selectors may be chosen measurably from the
nonempty compact set of solutions of these finite linear constraints;
for example take the unique minimum-Euclidean-norm feasible selector.
Zero training samples have zero gradient and may be assigned selector zero.
This argument obtains selectors from the Clarke inclusion; it does not
assert that every independently selected boundary backpropagation field
is a Clarke element.

Taking the parameter inner product of \eqref{eq:flow} gives
\[
 \frac{d}{dt}\|\theta\|^2
 =4\mathbb E_n\langle x-f_t(x),f_t(x)\rangle
 \le\mathbb E_n\|x\|^2.
\]
Consequently $\|\theta(t)\|^2\le\|\theta(0)\|^2+t\mathbb E_n\|x\|^2$.
On each finite time interval the vector field is bounded on this ball,
so local solutions continue globally. Polynomial maps and ReLU are
locally Lipschitz; their compositions with the trajectory are absolutely
continuous on compact time intervals. If $q$ is absolutely continuous,
$(q_+)'=\mathbf1_{\{q>0\}}q'$ almost everywhere: $q'=0$ almost everywhere
on $\{q=0\}$. This also justifies all probe-gate differentiations below.
Finally the Whitney-stratifiable chain rule applies to the semialgebraic
loss. Theorem~5.8 and Lemma~5.2 of \cite{davis2020} give
$\dot{\mathcal L}=\langle-\dot\theta,\dot\theta\rangle$,
which proves \eqref{eq:dissipation}.
\end{proof}

\begin{corollary}[Neuronwise balance]\label{cor:exact-balance}
If $\|w_i(0)\|=\|u_i(0)\|$ for every $i$, then
$\|w_i(t)\|=\|u_i(t)\|$ for every $i$ and every $t\ge0$.
\end{corollary}
\begin{proof}
Both squared norms have derivative
$2\mathbb E_n[\langle w_i,r_x\rangle(u_i^\top x)_+]$ by
\eqref{eq:flow}. Their difference is absolutely continuous and constant.
\end{proof}

\begin{lemma}[Mass and parameter-displacement envelope]\label{sbwindow:mass}
Put $C_x=\mathbb E_n\|x\|^2$ and
$\lambda=\sqrt{2\mathcal L(0)C_x}$. Under balanced initialization, for
any $\lambda_0\ge\lambda$ and $t\ge0$,
\begin{gather}\label{eq:mass-envelope}
 S(t)\le s e^{2\lambda_0t},\qquad
 \|U(t)-U(0)\|_F,\ \|W(t)-W(0)\|_F
 \le\sqrt{s}(e^{\lambda_0t}-1).
\end{gather}
For unit $\xi,v$, the change of $f_t(\xi)$ and of
$\sum_i(u_i^\top\xi)_+(u_i^\top v)_+$ is at most
$\Delta_{\rm prod}(t):=s(e^{2\lambda_0t}-1)$.
Moreover
\[
 (\|W\|_F+\|U\|_F)\|W-U\|_F
 \le D_*(t):=4s e^{\lambda_0t}(e^{\lambda_0t}-1).
\]
\end{lemma}
\begin{proof}
Cauchy--Schwarz in \eqref{eq:flow}, balance, and
\eqref{eq:dissipation} give
$\|\dot u_i\|,\|\dot w_i\|\le\lambda\|u_i\|$.
Thus $\dot S\le2\lambda S$ in physical time and each matrix speed is at most
$\lambda\sqrt S$. Integration proves \eqref{eq:mass-envelope}.
ReLU Lipschitzness gives each product displacement at most
$(\sqrt{s}e^{\lambda_0t}+\sqrt{s})\sqrt{s}(e^{\lambda_0t}-1)=\Delta_{\rm prod}(t)$.
The final bound follows from $W(0)=U(0)$ and the triangle inequality.
\end{proof}

\subsection{Exact instantaneous probe kernel}
\begin{lemma}[Probe kernel and residual pairing]\label{lem:probe-kernel}
For a fixed probe $\xi$, set $c_i(\xi)=\mathbf1_{\{u_i^\top\xi>0\}}$ and
\begin{equation}\label{sbaudit:kernel}
 \mathsf K_t(\xi,x)=\sum_i\big[(u_i^\top\xi)_+(u_i^\top x)_+I
 +c_i(\xi)\alpha_i(x)(\xi^\top x)w_iw_i^\top\big].
\end{equation}
Then almost everywhere,
\begin{equation}\label{sbaudit:alignment}
 \dot f_t(\xi)=\mathbb E_n[\mathsf K_t(\xi,x)r_x],\qquad
 \dot E(\xi,t)=-\mathbb E_n\langle r_\xi,\mathsf K_t(\xi,x)r_x\rangle.
\end{equation}
\end{lemma}
\begin{proof}
Differentiate the finite sum defining $f_t(\xi)$ using
Lemma~\ref{lem:gf-regularity}, substitute \eqref{eq:flow}, and interchange
the finite sums. The second equality is the squared-norm chain rule.
Probe selectors $c_i$ and training selectors $\alpha_i$ remain distinct;
no activation pattern is frozen.
\end{proof}


\section{Experimental protocol and additional empirical results}\label{app:empirical}
\subsection{Canonical experiment}
The settings and diagnostics below were checked against the archived
experimental code and result files. Identifying repository details are
withheld for anonymous review. The canonical
experiment is planar, with means $(3,0)$ and $(0,2)$, isotropic standard
deviation $0.15$, $2000$ samples per cluster ($N=4000$), width $h=200$,
and seed $0$. Each initialization row has norm $\varepsilon=10^{-3}$,
uniform angle, and $W(0)=U(0)$. The empirical gradient of
$\mathcal L=(2N)^{-1}\sum\|f(x)-x\|^2$ is evaluated with the actual ReLU
masks. Explicit Euler uses step $2\cdot10^{-4}$ through physical time
$10$; states are saved every $50$ steps, at spacing $0.01$.
These settings are defined in \path{full_flow_diagnostic.py} and
\path{theory_guided_swing_by_phase_diagram.py}.
They differ from the certified regime in Appendix~\ref{app:finite}.

\subsection{Directional phase diagram and a compositional probe}
For $x(\psi)=(\cos\psi,\sin\psi)$, the directional diagnostic computes
$E(x(\psi),t)$ from the saved full-network states. It smooths the error
with an eleven-sample moving average. For each local minimum at index $i$
and later index $j$, let $\Delta=\widetilde E_j-\widetilde E_i$.
The increase is persistent if, for the following $0.20$ time units,
$\widetilde E\ge\widetilde E_i+0.8\Delta$ at every saved time.
Define $S(\psi)$ as the largest positive persistent increase, or zero
when there is none. The reported reversal time is the minimizing index
for this selected increase. This sampled, smoothed statistic is distinct
from a zero of the continuous-time error derivative.

The phase diagram starts at the first logged weak-branch state and ends
at $10$. Its $1022$ directions combine a $720$-point global grid,
$0.05^\circ$ refinement near the swept weak-selector boundaries, and
explicit near-axis directions. The closed positive cone includes both
positive axes (a $10^{-12}$ coordinate tolerance handles trigonometric
roundoff): $198$ directions are inside and $824$ outside. Exactly $230$
directions have $S(\psi)>10^{-4}$, and all $230$ are outside the closed
positive cone. These are grid counts, not angular probabilities.

The zero in-cone count is threshold-dependent. The largest sampled
in-cone amplitude is $5.1392\cdot10^{-5}$, so the reported $10^{-4}$
cutoff lies about $1.95$ times above every sampled in-cone amplitude.
Lowering the cutoff to $5\cdot10^{-5}$ or $10^{-5}$ yields respectively
$2$ or $31$ in-cone directions; therefore the claim is not that the
in-cone and off-cone amplitude populations are disjoint. The data are in
\path{theory_guided_swing_phase_diagram_v3/phase_diagram.csv}.
Figure~\ref{fig:phase-diagram-app} summarizes this directional sweep and the
reported thresholding convention.

\begin{figure}[t]
\centering
\includegraphics[width=0.78\linewidth]{theory_guided_swing_phase_diagram_v3/phase_diagram_S_corrected.pdf}
\caption{Canonical directional phase diagram for the persistent-increase
statistic $S(\psi)$. The $10^{-4}$ reporting threshold excludes every
sampled direction in the closed positive cone, although smaller in-cone
increases appear at lower thresholds.}
\label{fig:phase-diagram-app}
\end{figure}

A separate complete-trajectory check uses the empirical positive-positive
probe
\[
 x^\star=(\widehat\mu_1,\widehat\mu_2)
 =(2.9957961621,\ 2.0009856079),
 \qquad
 \overline E(t):=\frac{E(x^\star,t)}{\|x^\star\|^2}.
\]
Its direction is $33.740196^\circ$, only $0.050129^\circ$ from the
population compositional direction $(3,2)/\sqrt{13}$ at
$33.690068^\circ$. Over the complete saved trajectory from $0$ to $10$,
rather than the branch-truncated window, the normalized sampled error
decreases from $0.4999513$ to $7.2666\cdot10^{-6}$; its largest
successive increment is $-2.66\cdot10^{-8}$. Thus no increase occurs at
the recorded times. The separate phase-diagram probe at exactly
$33.7^\circ$ likewise has $S=0$. These observations do not prove
continuous-time monotonicity.

\subsection{Strong-block decomposition}
At $t_{\rm block}=4$, freeze the strong and weak index sets within
$15^\circ$ of $e_1$ and the tracked weak root, respectively; the remaining
indices form the remainder set. The strong set has $71$ neurons.
The matrices are $M^{(p)}=W_{I_p}^\top U_{I_p}$, without probe gates.
For a fixed selector in the reduced model, each entry contributes
\begin{equation}\label{eq:empirical-entry}
 \dot E_{\rm red}=\sum_{r,s}(M v-v)_r\,\dot m_{rs}\,v_s.
\end{equation}
This is an entrywise decomposition of the reduced derivative, not of
the full-network derivative when the reduction residual is nonzero.
At the $-85^\circ$ probe the measured strong matrix at $t=4$ is
\[
 M^{(1)}=\begin{pmatrix}
 0.9343428&-0.0101619\\-0.0942215&0.0061059
 \end{pmatrix}.
\]
The stored entry contributions at the logged full reversal near $3.67$
are approximately
\[
\begin{array}{c|rrrr}
\text{entry}&m_{11}&m_{12}&m_{21}&m_{22}\\\hline
\text{contribution}&2.80\cdot10^{-6}&-1.47\cdot10^{-4}&-1.12\cdot10^{-3}&1.50\cdot10^{-3}
\end{array}
\]
The full persistent increase is $0.0090926$, versus $0.0100671$ for the
reduced model. The stored $0.20$-wide windows immediately before and after
the full reversal make the sign change more explicit:
\[
\begin{array}{c|rr}
\text{contribution}&\text{before}&\text{after}\\\hline
m_{11}& 3.78\cdot10^{-6}& 2.07\cdot10^{-6}\\
m_{12}&-1.54\cdot10^{-4}&-1.29\cdot10^{-4}\\
m_{21}&-1.298\cdot10^{-3}&-1.001\cdot10^{-3}\\
m_{22}& 1.274\cdot10^{-3}& 1.654\cdot10^{-3}
\end{array}
\]
Their summed reduced-derivative contribution changes from approximately
$-1.739\cdot10^{-4}$ before the reversal to
$+5.257\cdot10^{-4}$ after it. Thus the sign change is controlled by the
changing balance of non-major/transverse strong-block contributions, rather
than by the dominant $m_{11}$ response. Here the table reports
\emph{contributions}, not the raw matrix entries themselves.
Agreement of reversal timing does not imply exact amplitude prediction.
The source files are \path{decomp_m85.csv} and
\path{decomposition_summary.csv} in the same phase-diagram directory.
Figure~\ref{fig:m85-decomp-app} shows the full/reduced loss comparison and
the entrywise reduced-derivative terms.

\begin{figure}[t]
\centering
\begin{minipage}{0.48\linewidth}
\centering
\includegraphics[width=\linewidth]{theory_guided_swing_phase_diagram_v3/decomp_m85_loss.pdf}\\[-1mm]
{\small (a) Full and reduced OOD error.}
\end{minipage}\hfill
\begin{minipage}{0.48\linewidth}
\centering
\includegraphics[width=\linewidth]{theory_guided_swing_phase_diagram_v3/decomp_m85_terms.pdf}\\[-1mm]
{\small (b) Entrywise contributions to $\dot E_{\rm red}$.}
\end{minipage}
\caption{Canonical $-85^\circ$ strong-block decomposition. The strongest
reversal is reproduced by the reduced block, while the derivative sign
change is carried by transverse/non-major terms.}
\label{fig:m85-decomp-app}
\end{figure}

\subsection{Weak-selector analysis}
The tracked weak-root selector is
$g_W(x,t)=\mathbf1_{\{\cos\phi_W(t)x_1+\sin\phi_W(t)x_2>0\}}$.
A switch is a sign change of this score between successive logged times;
the crossing time is linearly interpolated. For a switch, with reduced
strong background $e_0=g_1M^{(1)}x-x$ and $q=M^{(2)}x$, the predicted
change on activation is $\langle e_0,q\rangle+\|q\|^2/2$ and its negative
on deactivation. This is a discrete selector-model comparison; the
full ReLU output itself is continuous.
The $26$ largest-amplitude reversals overall, spanning approximately
$-88.5^\circ$ to $-76^\circ$, undergo no weak-selector switch. Thus a
selector switch is not necessary for the strongest observed reversals.

Separately, restrict to directions that undergo a selector switch, satisfy
the reporting threshold $S(\psi)>10^{-4}$, and have a finite observed
reversal time. There are $126$ such directions. Ranking this
\emph{switching, thresholded} set by observed persistent increase and
examining the switch nearest the observed reversal, the stored signed
selector prediction is positive in all $126$ cases. Without the reversal
threshold, $167$ of the $205$ switching directions with a defined signed
prediction are positive. The thresholded selector analysis is
\emph{post hoc}: it is not a preregistered test of reversal timing and is
not used as a theorem premise.

\subsection{Nine-cell protocol and timing results}
The preregistered grid fixes $\mu_1=3$, takes
$\mu_2\in\{1.8,2.0,2.2\}$ and $\sigma\in\{0.12,0.15,0.18\}$,
and uses seed $0$ in every cell. Its canonical and grid results therefore
use a single seed. Before evaluating full-ReLU OOD outcomes, the pilot
protocol was amended to prevent eventual block members from being
extrapolated backward before recruitment. For each frozen block, let
$t_{W,p}$ be the first time its projected-mass share stays at least
$90\%$ of its $t_{\rm block}$ share through $t_{\rm block}$.
The amended validity start is
$t_W=\max\{t_{\rm birth},t_{W,1},t_{W,2}\}$, and predictions are evaluated
on $[t_W,10]$. Existing trajectories were reused without retraining.
The correction record explicitly states
\texttt{full\_ood\_seen\_before\_correction: false}:
full OOD outcomes had not been inspected before the correction.
The corrected protocol, per-cell predictions and SHA-256 manifest were
frozen before the separate full-output evaluation.

The choice $T_{\rm BLOCK}=4$ is \emph{offline localization}, not a
forecasting horizon. Block membership and $t_W$ use the in-distribution
trajectory up to that time; they are not prospective predictions at $t=0$.
They were not selected using future full-network OOD reversal times.
The operational scale, read from the block at time $4$, is
\[
 \eta_p^{\rm blk}=\max\left\{
 \frac{|m_{pq}|}{|m_{pp}|},\frac{|m_{qp}|}{|m_{pp}|},
 \sqrt{\frac{|m_{qq}|}{|m_{pp}|}}\right\}.
\]
Core-profile tests use $|z|\le4$; the timing probe is specifically
$z=1$, $s_\xi=-1$. It is not selected from the full OOD peak.
The recorded reduced core peak is a separate statistic.
The grid uses finite differences of the saved block matrices to compute
\eqref{eq:eta-independent}. Negative-to-positive crossings are linearly
interpolated; the rule searches within one time unit of the reduced-model
reversal, with the nearest valid crossing as a fallback. Every cell had
exactly one $G$ crossing after $t_W$ (and in its relevant search window),
so the rule never had to choose among multiple crossings.

\begin{center}
\begin{tabular}{ccrrr}\toprule
$\mu_2$&$\sigma$&$t_G$&$t_{\rm rev,red}$&$t_{\rm rev,full}$\\\midrule
1.8&0.12&4.4294&4.42&4.52\\
1.8&0.15&4.3643&4.36&4.44\\
1.8&0.18&4.3088&4.31&4.36\\
2.0&0.12&3.7200&3.73&3.79\\
2.0&0.15&3.6479&3.66&3.73\\
2.0&0.18&3.5914&3.58&3.65\\
2.2&0.12&3.2094&3.22&3.24\\
2.2&0.15&3.1041&3.11&3.16\\
2.2&0.18&3.0793&3.08&3.12\\\bottomrule
\end{tabular}
\end{center}
These entries come from
\path{swing_grid_preregistered/evaluation_v1_postspecialization/G_timing_summary.csv}.
Mean absolute timing errors are $0.0073$ against the reduced reversal and
$0.0617$ against the full reversal. The nine-cell grid remains an empirical
study rather than an unconditional sign theorem; at the canonical setting,
Appendix~\ref{app:canonical-cert} gives a separate conditional continuous-time
certificate for the sign ordering itself. The stored scaling study supports
order-one $S_{\max}/(\eta_p^{\rm blk})^2$ over this grid, but not a universal
coefficient. The weak-side profiles do not collapse into a single common
curve. The parameter-only prediction $\eta_1\propto\mu_2\sigma/\mu_1^2$
and the fixed-gate equilibrium proxy are not supported by this grid;
measured block coefficients are retained rather than replacing them by
those predictions.

\subsection{Canonical residual diagnostics and separate robustness audit}
For the frozen strong block and $-85^\circ$ probe, evaluate the empirical
vector field on the saved states in $[2.5,5]$, spaced by $0.01$.
The following are maxima or correlations over those recorded states,
not continuous-time certified bounds:
\begin{center}
\begin{tabular}{lr}\toprule
Diagnostic&Measured value\\\midrule
$\max\|R\|$&$5.7908\cdot10^{-3}$\\
$\max\|\dot R\|$&$4.1682\cdot10^{-3}$\\
$\max\|R\|/\|f(v)-v\|$&$0.005880$\\
$\max|\eta^2\mathcal G-\dot E_{\rm red}|$&$1.8401\cdot10^{-4}$\\
Correlation of $\eta^2\mathcal G$ and $\dot E_{\rm red}$&$0.9999985$\\
Interpolated zeros: $G$, reduced derivative, full derivative&$3.6018,\ 3.6035,\ 3.6726$\\\bottomrule
\end{tabular}
\end{center}
The analytic derivatives use the actual training masks at each state.
There are no active weak-block neurons. Exactly five of the $71$
strong-block neurons are inactive at every one of the $251$ sampled
states; their ungated contribution accounts for almost all of $R$.
Exactly $38$ remainder neurons are active at every sampled state, and
their combined output norm is at most $2.135\cdot10^{-5}$. Thus small
reduction error here does not mean all strong selectors are active.
These quantities were recomputed from the saved trajectory, without a
new training run, and are supplied in the accompanying
\path{CANONICAL_RESIDUAL_DIAGNOSTIC.csv}. The existing canonical $G$
script alone does not output these residual and active-count diagnostics.
Figure~\ref{fig:canonical-G-app} displays the corresponding canonical timing
trace.

\begin{figure}[t]
\centering
\includegraphics[width=0.72\linewidth]{canonical_G_check/canonical_G_edot.pdf}
\caption{Canonical timing diagnostic at the $-85^\circ$ probe. The scaled
normal-form quantity $\eta^2\mathcal G$ tracks the reduced-model error
derivative, while the full-network zero occurs slightly later.}
\label{fig:canonical-G-app}
\end{figure}

The archived results also contain a \emph{separate} robustness audit:
three cells $(1.8,0.12)$, $(2.0,0.15)$ and $(2.2,0.18)$, each at seeds
$0,1,2,3,4$ ($15$ runs, reusing seed $0$). Its frozen protocol and
\path{swing_final_audits/final_audit_summary.json} report mean absolute
$G$ timing error $0.04897$ and maximum $0.09058$. This additional study
does not change the seed count or preregistration history of the nine-cell
table. Figure~\ref{fig:multiseed-G-app} summarizes this separate robustness
audit.

\begin{figure}[t]
\centering
\includegraphics[width=0.72\linewidth]{swing_final_audits/multiseed/multiseed_G_timing_error.pdf}
\caption{Separate five-seed robustness audit on three representative grid
cells. This audit is not part of the nine-cell preregistered test.}
\label{fig:multiseed-G-app}
\end{figure}


\section{Pointwise OOD normal form}\label{app:normal}
\subsection{Frozen block and exact remainder}
Fix a set $I_{\rm str}\subseteq[h]$ throughout the interval of interest and
use its \emph{ungated} matrix
\[
 M_{\rm str}(t)=W_{I_{\rm str}}(t)^\top U_{I_{\rm str}}(t),\qquad
 R(t;v)=f_t(v)-M_{\rm str}(t)v.
\]
The identity is exact. The remainder includes inactive strong neurons,
weak neurons and unassigned neurons; it is not just unassigned mass.
The following result concerns a planar block in the ordered basis
$(e_p,e_q)$, $p\ne q$. It does not assert that training produces its
coefficient bounds or a small remainder.

\begin{proposition}[Pointwise block normal form]\label{prop:pointwise-normal}
Fix $0<\eta\le1$, $s_\xi\in\{-1,1\}$ and $|z|\le Z$ with $\eta|z|<1$.
Let $k=\sqrt{1-\eta^2z^2}$ and $v=\eta z e_p+s_\xi k e_q$.
Suppose on a compact time interval that
\[
 M_{\rm str}=\begin{pmatrix}a&\eta b\\\eta c&\eta^2d\end{pmatrix}
\]
has absolutely continuous coefficients with
$|a|,|b|,|c|,|d|,|\dot a|,|\dot b|,|\dot c|,|\dot d|\le K$
(the derivative bounds hold almost everywhere). Suppose
$\|R(t;v)\|\le\epsilon_{\rm red}\le1$ and
$\|\dot R(t;v)\|\le\epsilon_{\rm dyn}\le1$ almost everywhere.
Define
\begin{align}
 \mathcal H&=\tfrac12[(a-1)z+s_\xi b]^2-s_\xi cz-d-\tfrac12z^2,
 \label{eq:normal-H}\\
 \mathcal G&=[(a-1)z+s_\xi b][\dot a z+s_\xi\dot b]
              -s_\xi z\dot c-\dot d.\label{eq:normal-G}
\end{align}
Uniformly in this interval,
\begin{align*}
 E(v,t)&=\tfrac12+\eta^2\mathcal H+\mathcal R_E,
 &|\mathcal R_E|&\le C_{K,Z}(\eta^4+\epsilon_{\rm red}),\\
 \dot E(v,t)&=\eta^2\mathcal G+\mathcal R_{\dot E},
 &|\mathcal R_{\dot E}|&\le C_{K,Z}(\eta^4+\epsilon_{\rm red}+\epsilon_{\rm dyn})
 \quad\text{a.e.}
\end{align*}
\end{proposition}
\begin{proof}
Set $A_k=(a-1)z+s_\xi k b$ and $T_k=cz+s_\xi k d$.
For $E_0=\tfrac12\|M_{\rm str}v-v\|^2$, direct expansion gives
\begin{align}\label{eq:normal-exact}
 E_0&=\tfrac12+\eta^2\{\tfrac12A_k^2-s_\xi kT_k-\tfrac12z^2\}
                     +\tfrac12\eta^4T_k^2,\\
 \dot E_0&=\eta^2\{A_k\dot A_k+(\eta^2T_k-s_\xi k)\dot T_k\}.
 \nonumber
\end{align}
Since $|1-k|=\eta^2z^2/(1+k)\le\eta^2Z^2$, replacing $k$ by one
in the braces costs $O_{K,Z}(\eta^2)$. All displayed coefficients and
rates are bounded by constants depending only on $K,Z$; hence both
errors in \eqref{eq:normal-exact} are $O_{K,Z}(\eta^4)$.
The exact remaining identities are
\begin{align*}
 E-E_0&=\langle M_{\rm str}v-v,R\rangle+\tfrac12\|R\|^2,\\
 \dot E-\dot E_0&=\langle\dot M_{\rm str}v,R\rangle
 +\langle M_{\rm str}v-v,\dot R\rangle+\langle R,\dot R\rangle.
\end{align*}
Because $\eta\le1$, the block norm and its derivative are bounded by
constants depending only on $K$. Cauchy--Schwarz proves the residual
bounds. Absolute continuity justifies the identities through probe-gate
changes without differentiating a gate indicator.
\end{proof}

\begin{remark}[When $\mathcal G$ determines the derivative sign]
The expansion is informative at the sign level only when the leading
normal-form margin dominates its remainder. A sufficient pointwise
condition is
\[
 \eta^2|\mathcal G(t)|
 >
 C_{K,Z}\bigl(\eta^4+\epsilon_{\rm red}+\epsilon_{\rm dyn}\bigr).
\]
Under this inequality, $\dot E(v,t)$ has the sign of $\mathcal G(t)$.
The empirical residual diagnostics below measure the relevant discrepancy
at the canonical probe, but they are not uniform analytic bounds proving
this condition from random initialization. Appendix~\ref{app:canonical-cert}
provides a separate computer-assisted conditional theorem at the canonical
parameters: once the flow enters the certified learned-state neighborhood,
$Q_G=\eta^2\mathcal G$ is rigorously positive before the full OOD error stops
decreasing. This conclusion is still independent of the finite-sample theorem
in Appendix~\ref{app:finite} and does not certify initialization-to-entry.
\end{remark}

\subsection{The empirical timing statistic}
Let $x_p=\eta z$ and let $m_{rs}$ denote the unscaled block entries.
Substitution into \eqref{eq:normal-G} gives exactly
\begin{equation}\label{eq:eta-independent}
 \eta^2\mathcal G=
 [(m_{pp}-1)x_p+s_\xi m_{pq}]
 [\dot m_{pp}x_p+s_\xi\dot m_{pq}]
 -s_\xi x_p\dot m_{qp}-\dot m_{qq}.
\end{equation}
For a fixed probe and fixed block, this expression is independent of the
operational choice of $\eta$. Different choices change the scaled
coefficients and $z$, and therefore whether the bounded-coefficient
hypotheses provide a useful approximation; they do not tune the value
or zero crossing of \eqref{eq:eta-independent}.
The canonical diagnostic differentiates the block analytically,
$\dot M=\dot W_I^\top U_I+W_I^\top\dot U_I$, using the empirical
vector field. The grid diagnostic uses finite differences of saved block
matrices. These two derivative conventions are reported separately.

\subsection{Empirical reduction diagnostics}
The canonical $-85^\circ$ probe is evaluated on $[2.5,5]$ at spacing
$0.01$. Appendix~\ref{app:empirical} reports measured $R,\dot R$,
active counts and derivative discrepancies. These are numerical diagnostics,
not uniform analytic bounds supplying $\epsilon_{\rm red}$ or
$\epsilon_{\rm dyn}$. In particular, the strong block has inactive neurons
at this probe. Its ungated reconstruction remains the intended interface;
those activation mismatches are included in $R$.

In the noiseless model, Proposition~\ref{sbsector:profile} proves
$R=0$, the required coefficient bounds, and explicit uniform error and
derivative estimates along the actual trajectory. It provides a proved
special case complementary to the empirical measurements in
Appendix~\ref{app:empirical}.



\section{Computer-assisted canonical continuation and timing}
\label{app:canonical-cert}

This section concerns the exact supplied canonical finite dataset and the
archived canonical reference trajectory. It is a deterministic,
computer-assisted statement for that dataset, not a high-probability theorem.
The canonical parameters are
\[
 (\mu_1,\mu_2,\sigma)=(3,2,0.15),\qquad N=4000,\qquad h=200,
\]
with the fixed $71$-neuron strong cohort selected by the archived
in-distribution rule at time $4$. The theorem applies to every compatible
Clarke flow satisfying the stated entry condition. No global uniqueness of
the nonsmooth flow and no persistent training-gate assumption are used.

\paragraph{Arithmetic and input scope.}
The verifier uses the supplied binary dataset and archived parameter states.
It stores midpoint--radius intervals. Basic floating-point sums and dot
products are evaluated in IEEE-754 round-to-nearest and enclosed with explicit
$\gamma_k=ku/(1-ku)$ error factors (with unit roundoff $u$), together with
outward \texttt{nextafter} inflation of scalar interval endpoints; interval
trigonometric endpoints are evaluated at high precision. The $\gamma_k$
bounds depend only on an upper bound for the number of rounded operations and
therefore cover any summation order using no more than that many operations.
Unsafe reassociation and fast-math are excluded. The supplied binary dataset
matches the pinned canonical generator under the recorded runtime, but the
historical trajectory archive did not itself preserve an independent raw-data
hash. Accordingly, the theorem is stated for the exact supplied finite
dataset rather than as an iid-data event.

Put
\[
 r=\sin(\pi/36),\qquad k=\cos(\pi/36),\qquad
 \xi_\alpha=(\cos\alpha,\sin\alpha),
\]
and define the explicit OOD sector
\begin{equation}\label{eq:canonical-sector}
 \mathcal I_{\rm can}
 :=\{\alpha:|\alpha+17\pi/36|<10^{-6}\}.
\end{equation}
Thus the central direction is $-85^\circ$ and the sector has width
$2\cdot10^{-6}$ radians. Let
\[
 M_I=\sum_{i\in I}w_i u_i^\top=(m_{ab})_{a,b=1}^2
\]
be the fixed ungated strong-block matrix. The repository timing statistic is
\begin{equation}\label{eq:canonical-QG}
 Q_G=[(m_{11}-1)r-m_{12}](r\dot m_{11}-\dot m_{12})
        +r\dot m_{21}-\dot m_{22}.
\end{equation}
By \eqref{eq:eta-independent}, this is exactly the canonical value of
$\eta^2\mathcal G$ for the $z=1$, $s_\xi=-1$ timing probe. The factors $k$
are not inserted into this auxiliary statistic.


\subsection{Enclosure inequalities used by the computer-assisted proofs}
\label{app:canonical-enclosures}

For completeness, we state the deterministic inequalities evaluated by the
canonical verifier. The code supplies outward numerical bounds for the
quantities below; the inequalities themselves are ordinary real-arithmetic
statements. Throughout this subsection use the residual convention
\[
 e_\theta(x):=x-f_\theta(x),\qquad
 F(\theta):=J_\theta^*e_\theta,
\]
so that $F$ is a compatible gradient-flow field. At a zero training
preactivation, all compatible selectors in $[0,1]$ are allowed.

\paragraph{Static neighborhoods.}
Fix a parameter center $y$, a radius $R>0$, and write
\[
 P=\|y\|,\qquad \bar P=P+R,\qquad D=PR+\frac12R^2.
\]
For a cluster index set $C$, use the full-dataset normalization
\[
 \|z\|_C^2:=\frac1N\sum_{n\in C}\|z_n\|^2,
 \qquad
 \lambda_C\ge
 \left\|\frac1N\sum_{n\in C}x_nx_n^\top\right\|_{\rm op}.
\]
Let $F_C$ be cluster $C$'s contribution to the full field, always using the
\emph{full-network} residual, and let $F_C^0$ be its polynomial extension
with the training masks frozen at $y$. For either potential used below, put
$g=\nabla\Phi$, $G_0=\|g(y)\|$. Let $v$ denote the fixed input vector used by
that potential, $I$ its fixed neuron index set, $P_I=\|(u_i,w_i)_{i\in I}(y)\|$,
and $\ell_y$ its output-side gradient. Define
\begin{align}
 L_C&:=\lambda_C\bar P^2+
 \sqrt{\lambda_C}\bigl(\|e_y\|_C+\sqrt{\lambda_C}D\bigr),
 \label{eq:canonical-LC}\\
 H_g&:=\|v\|^2(P_I+R)^2+
 \|v\|\Bigl(\|\ell_y\|+\|v\|(P_IR+R^2/2)\Bigr).
 \label{eq:canonical-Hg}
\end{align}
To enclose changes of training gates, set
\[
 q_{ni}=u_i(y)^\top x_n,
 \qquad s_n=\|x_n\|,
 \qquad e_n=e_y(x_n),
\]
\begin{align*}
 d_{ni}&=(Rs_n-|q_{ni}|)_+,
 &p_{ni}&=\mathbf1_{\{|q_{ni}|\le Rs_n\}},\\
 b_{ni}&=|w_i(y)^\top e_n|+R\|e_n\|
       +(\|w_i(y)\|+R)Ds_n,\\
 j_i^u&=\frac1N\sum_{n\in C}p_{ni}b_{ni}s_n,
 &j_i^w&=\frac1N\sum_{n\in C}d_{ni}(\|e_n\|+Ds_n),\\
 z_n&=\sum_i(\|w_i(y)\|+R)d_{ni},
\end{align*}
and
\begin{equation}
 B_C:=\left[\sum_i\bigl((j_i^u)^2+(j_i^w)^2\bigr)\right]^{1/2}
       +\sqrt{\lambda_C}\,\bar P\,\|z\|_C.
 \label{eq:canonical-BC}
\end{equation}

\begin{proposition}[Static neighborhood enclosure]
\label{prop:canonical-static-neighborhood}
For every compatible training-selector realization and every
$\|\theta-y\|\le R$,
\begin{align}
&\left|g(\theta)^TF_C(\theta)-g(y)^TF_C(y)\right|\nonumber\\
&\quad\le
 R\left[H_g\bigl(\|F_C(y)\|+L_CR\bigr)+G_0L_C\right]
 +(G_0+H_gR)B_C.
 \label{eq:canonical-static-neighborhood}
\end{align}
The same inequality applies to the total derivative and separately to each
cluster's signed contribution.
\end{proposition}

\begin{proof}
Cauchy--Schwarz over neuron and input/output parameter coordinates gives
$\|J_C\|\le\sqrt{\lambda_C}\|\theta\|$. Integration along the parameter
segment gives
\[
 \|f_\theta-f_y\|_C\le\sqrt{\lambda_C}D,
 \qquad
 \|f_\theta(x_n)-f_y(x_n)\|\le Ds_n,
\]
and the same bounds hold for the fixed-mask polynomial extension. On a
smooth piece,
$DF_C^0=-J_C^*J_C+H_{e,C}$, and the neuronwise residual Hessian has norm at
most $\sqrt{\lambda_C}\|e\|_C$, so $\|DF_C^0\|\le L_C$. Direct
differentiation of either potential gives $\|Dg\|\le H_g$ on the ball.
Adding and subtracting $g(y)^TF_C^0(\theta)$ yields the first term of
\eqref{eq:canonical-static-neighborhood}.

A changed activation differs from its fixed-mask extension by at most
$d_{ni}$ and its selector can differ only when $p_{ni}=1$. The pointwise
residual and weight bounds imply
$|w_i(\theta)^Te_\theta(x_n)|\le b_{ni}$, hence the direct input- and
output-gradient changes are bounded by $j_i^u$ and $j_i^w$. The residual
change between the actual network and the fixed-mask extension is bounded
pointwise by $z_n$; applying the fixed-mask adjoint to this residual costs
at most $\sqrt{\lambda_C}\bar P\|z\|_C$. Thus
$\|F_C-F_C^0\|\le B_C$. Finally
$\|g(\theta)\|\le G_0+H_gR$, proving the claim. Compatible partial selectors
at boundary points obey the same bounds.
\end{proof}

\paragraph{One-sided tube propagation through gate jumps.}
Fix a reference center $y$ and radius $R$. On the smooth activation pieces
intersecting the ball, suppose the residual-Hessian contribution has
operator norm at most $a$. For each neuron/sample pair choose a uniform
upper bound
\[
 h_{in}\ge \frac1N\|x_n\|\,(w_i^Te_\theta(x_n))_+
\]
throughout the ball, and for nonzero samples define the normalized reference
margin
\[
 d_{in}:=\frac{|u_i(y)^Tx_n|}{\|x_n\|},
 \qquad
 B_i(s):=\sum_{n:d_{in}\le s}h_{in},\quad 0\le s\le R.
\]
If the center ranges over a reference subinterval, uniform lower bounds on
$d_{in}$ and uniform upper bounds on $h_{in}$ may be used instead.

\begin{lemma}[One-sided stability with training-gate jumps]
\label{lem:canonical-tube-stability}
Suppose
$B_i(s)\le a_i+\beta s$ for every $i$ and every $0\le s\le R$, with
$a_i,\beta\ge0$. Then, for $\|z\|\le R$,
\begin{equation}
 \langle z,F(y+z)-F(y)\rangle
 \le(a+\beta)\|z\|^2+\|(a_i)_i\|_2\|z\|.
 \label{eq:canonical-one-sided-stability}
\end{equation}
Compatible endpoint selectors and simultaneous crossings are included. If
$y(t)$ is an absolutely continuous reference satisfying
$\|F(y(t))-\dot y(t)\|\le\eta$ and
$r(t)=\|\theta(t)-y(t)\|$, then while $r\le R$,
\begin{equation}
 D^+r\le(a+\beta)r+\|(a_i)_i\|_2+\eta
 \label{eq:canonical-radius-ode}
\end{equation}
in the scalar comparison sense.
\end{lemma}

\begin{proof}
On each smooth activation cell,
$DF=-J^*J+H_e$; the first term is nonpositive in the quadratic pairing and
the second contributes at most $a\|z\|^2$. Integrate this bound along the
segment $y+sz$ and separate the finitely many training-hyperplane crossings.
At a crossing of sample $n$ for neuron $i$, the vector-field jump has only an
input-coordinate component, because the ReLU activation itself is zero there.
Its pairing with $z$ is at most
$h_{in}\|z_{u_i}\|$. Such a crossing is possible only if
$d_{in}\le\|z_{u_i}\|$, hence the total jump contribution is at most
\[
 \sum_i\|z_{u_i}\|B_i(\|z_{u_i}\|)
 \le \|(a_i)_i\|_2\|z\|+\beta\|z\|^2.
\]
Partial jumps at endpoints obey the same inequality, and segments lying in a
training hyperplane have zero conormal pairing for that gate. This proves
\eqref{eq:canonical-one-sided-stability}. Applying the bound to the error
equation $\theta-y$ and adding the reference defect gives
\eqref{eq:canonical-radius-ode}; standard norm regularization gives the
comparison at $r=0$.
\end{proof}

For the finite set of sample distances used by the verifier, a permissible
intercept is explicit:
\begin{equation}
 a_i(\beta)=\max\left\{0,\max_{0\le s\le R}(B_i(s)-\beta s)\right\},
 \label{eq:canonical-intercept}
\end{equation}
and the inner maximum needs only the sorted distances, including zero.

\begin{corollary}[Verified step criterion]
\label{cor:canonical-verified-step}
On a reference interval of length $h$, assume the preceding bounds hold
uniformly on its radius-$R$ tube. Put
$k=a+\beta$ and $\omega=\|(a_i)_i\|_2+\eta$. If the initial error obeys
$r_0<R$ and
\begin{equation}
 e^{kh}r_0+\omega\frac{e^{kh}-1}{k}<R,
 \label{eq:canonical-step-criterion}
\end{equation}
then the flow remains in the tube on the whole interval, and its terminal
error is bounded by the left-hand side. For $k=0$, the quotient in
\eqref{eq:canonical-step-criterion} is interpreted as $h$.
\end{corollary}

\begin{proof}
Apply the scalar integrating-factor inequality to
\eqref{eq:canonical-radius-ode} until a hypothetical first exit. Its upper
bound is nondecreasing in time and, by
\eqref{eq:canonical-step-criterion}, remains strictly below $R$ at the
terminal time. A first exit is therefore impossible. Iterating the terminal
bound gives the multistep enclosure.
\end{proof}

\paragraph{Cubic-Hermite reference defect.}
The local and timing verifiers use cubic Hermite reference segments of exact
length $h=1/5000$. Center one such step at $s=0$ and put
$a_h=h/2=1/10000$. Freeze the training masks at the midpoint and denote the
resulting polynomial field by $F^0$. Let
\begin{align*}
 d_0&=\|F^0(y)-y'\|,\\
 d_1&=\|DF^0(y)y'-y''\|,\\
 d_2&=\|D^2F^0(y)[y',y']+DF^0(y)y''-y'''\|.
\end{align*}
Suppose $P_*,V_*,A_*,T_*$ bound respectively
$\|y\|,\|y'\|,\|y''\|,\|y'''\|$ throughout the step and
$L_*$ bounds $\|DF^0\|$. If
$\lambda\ge\|N^{-1}\sum_nx_nx_n^T\|_{\rm op}$, then
\[
 \|D^2F^0\|\le3\lambda P_*,\qquad
 \|D^3F^0\|\le3\lambda.
\]
Because the reference is cubic, $y''''=0$, and Taylor's integral remainder
gives the uniform fixed-mask defect bound
\begin{equation}
 \eta_{\rm sm}:=
 d_0+a_hd_1+\frac12a_h^2d_2+
 \frac16a_h^3\left(
 3\lambda V_*^3+9\lambda P_*V_*A_*+L_*T_*
 \right).
 \label{eq:canonical-smooth-defect}
\end{equation}
The actual reference defect must also include training gates crossed by the
reference polynomial. For each neuron, its input excursion from the midpoint
is bounded by
\begin{equation}
 d_i^u\le a_h\|y'_{u_i}(0)\|+\frac12a_h^2\|y''_{u_i}(0)\|
       +\frac16a_h^3\|y'''_{u_i}\|,
 \label{eq:canonical-input-excursion}
\end{equation}
with the analogous $d_i^w$. Let
$m_{ni}=|u_i(y)^Tx_n|/\|x_n\|$, let
$p_{ni}^{\rm ref}=\mathbf1_{\{m_{ni}\le d_i^u\}}$, and put
\[
 a_{ni}^{\rm ref}:=(d_i^u\|x_n\|-|u_i(y)^Tx_n|)_+.
\]
If $\Delta_{\rm ref}$ is the whole-parameter Hermite excursion and
$D_{\rm ref}=P\Delta_{\rm ref}+\Delta_{\rm ref}^2/2$, define
\begin{align*}
 \epsilon_n^{\rm ref}&:=\|e_y(x_n)\|+D_{\rm ref}\|x_n\|,\\
 r_{ni}^{\rm ref}&:=|w_i(y)^Te_y(x_n)|
  +\|e_y(x_n)\|d_i^w
  +D_{\rm ref}\|x_n\|(\|w_i(y)\|+d_i^w),\\
 j_i^{u,{\rm ref}}&:=\frac1N\sum_n
  p_{ni}^{\rm ref}r_{ni}^{\rm ref}\|x_n\|,\\
 j_i^{w,{\rm ref}}&:=\frac1N\sum_n
  a_{ni}^{\rm ref}\epsilon_n^{\rm ref},\\
 z_n^{\rm ref}&:=\sum_i a_{ni}^{\rm ref}(\|w_i(y)\|+d_i^w).
\end{align*}
With $P_{\rm ref}=P+\Delta_{\rm ref}$, the reference-path gate defect used
by the verifier is bounded by
\begin{equation}
 \eta_{\rm gate}\le
 \left[\sum_i\bigl((j_i^{u,{\rm ref}})^2+(j_i^{w,{\rm ref}})^2\bigr)\right]^{1/2}
 +\sqrt\lambda\,P_{\rm ref}
   \left(\frac1N\sum_n(z_n^{\rm ref})^2\right)^{1/2}.
 \label{eq:canonical-gate-defect}
\end{equation}
Thus the certified defect is
\begin{equation}
 \eta=\eta_{\rm sm}+\eta_{\rm gate}.
 \label{eq:canonical-total-defect}
\end{equation}
This gate term is essential: the favorable sign in
Lemma~\ref{lem:canonical-tube-stability} concerns stability of two nearby
states and cannot be reused to omit gate changes of the numerical reference
itself.

For the one-sided rate on a radius-$R$ tube, the implementation uses the
midpoint residual matrices
$A_i^0=N^{-1}\sum_n\alpha_{ni}e_y(x_n)x_n^T$, the total output-change coefficient
$D_{\rm tot}$ on the reference excursion plus the tube, and flags
$p_{ni}^{\rm tube}$ for every potentially changed training gate. A valid
smooth-piece coefficient is
\begin{equation}
 a_*:=\max_i\|A_i^0\|+\lambda D_{\rm tot}
 +\max_i\frac1N\sum_n p_{ni}^{\rm tube}
       \|e_y(x_n)\|\|x_n\|.
 \label{eq:canonical-rate-instantiation}
\end{equation}
The positive conormal weights are bounded on that same tube, sorted by their
outward lower bounds on the reference-path gate distances, and inserted into
\eqref{eq:canonical-intercept}. The wide canonical run takes $\beta=1$ and
checks Corollary~\ref{cor:canonical-verified-step} with these quantities on
every step. The later sharper gate estimates developed in the research audit
are not premises of Theorems~\ref{thm:canonical-local-continuation}
and~\ref{thm:canonical-full-tube-timing}.

\subsection{Certified continuation from a learned-state neighborhood}

Let $\widehat\theta_{3.4}$ denote the exact supplied numerical reference
state at time $3.4$.

\begin{theorem}[Canonical continuation from a learned-state neighborhood]
\label{thm:canonical-local-continuation}
If a global compatible Clarke flow satisfies
\begin{equation}\label{eq:canonical-entry}
 \|\theta(3.4)-\widehat\theta_{3.4}\|_2
 \le1.2\cdot10^{-5},
\end{equation}
then, uniformly for every $\alpha\in\mathcal I_{\rm can}$,
\begin{align}
 E(\xi_\alpha,3.4)-E(\xi_\alpha,3.7)&>1.2\cdot10^{-4},
 \label{eq:canonical-certified-drop}\\
 E(\xi_\alpha,3.9)-E(\xi_\alpha,3.7)&>4\cdot10^{-5}.
 \label{eq:canonical-certified-rebound}
\end{align}
Consequently every such probe has an interior minimizer on $[3.4,3.9]$.
At entry the fixed strong cohort satisfies
\[
 m_{11}>0.9323,
\]
its input directions lie within $15^\circ$ of the strong axis and its output
directions within $20^\circ$. For the central probe, both $Q_G$ and the full
OOD rate are negative at entry and positive at $3.9$ for every compatible
training-selector realization. The strong-cluster contribution to the full
OOD rate is negative at both endpoints, the weak-cluster contribution is
positive, and the weak contribution dominates at the late endpoint.
Moreover, almost everywhere and uniformly on the sector,
\[
 \dot E<-0.0012\quad\text{on }[3.4,3.40001],
 \qquad
 \dot E>0.00040\quad\text{on }[3.9,3.90001].
\]
No initialized entry guarantee and no unique continuous zero of $Q_G$ is
asserted.
\end{theorem}

\begin{proof}
The continuation certificate uses $2500$ reference intervals of exact length
$1/5000$ on $[3.4,3.9]$. On each interval,
Lemma~\ref{lem:canonical-tube-stability}, the cubic defect bound
\eqref{eq:canonical-total-defect}, and
Corollary~\ref{cor:canonical-verified-step} enclose every compatible training
selector, including selectors at zero preactivation. An independent scalar
replay starts from the upward enclosure of the exact decimal radius
$1.2\cdot10^{-5}$ and checks every strict step inequality
\eqref{eq:canonical-step-criterion}. The certified trajectory radii at the two later snapshots obey
\[
 \|\theta(3.7)-\widehat\theta_{3.7}\|<2.828\cdot10^{-5},
 \qquad
 \|\theta(3.9)-\widehat\theta_{3.9}\|<4.251\cdot10^{-5}.
\]

Proposition~\ref{prop:canonical-static-neighborhood} is then applied on
radii $2\cdot10^{-5}$, $3.5\cdot10^{-5}$ and $4.6\cdot10^{-5}$ at times
$3.4,3.7,3.9$.
For the central probe $\xi_0=\xi_{-17\pi/36}$ they give
\begin{align*}
 E(\xi_0,3.4)&\in[0.48314641,0.48315250],\\
 E(\xi_0,3.7)&\in[0.48298853,0.48299930],\\
 E(\xi_0,3.9)&\in[0.48305831,0.48307257].
\end{align*}
The full-precision central gaps exceed $1.4712\cdot10^{-4}$ and
$5.902\cdot10^{-5}$. The largest certified full-parameter norm over the three
static balls is less than $2.006919462$, so
\[
 L:=\sum_i\|w_i\|\|u_i\|\le\frac12\|\theta\|^2
 <2.013862863,
 \qquad (1+L)^2<9.083369357<9.084.
\]
Proposition~\ref{sbaudit:secants} therefore gives the required unit-sphere
error Lipschitz constant. Transport over angular distance below $10^{-6}$
costs less than $1.8168\cdot10^{-5}$ for a difference of two errors. This leaves the sector-wide margins
\eqref{eq:canonical-certified-drop}--\eqref{eq:canonical-certified-rebound}.

The endpoint interval bounds for $10^3$ times the exact rates are
\begin{center}
\begin{tabular}{llrr}
\toprule
Statistic & contribution & $t=3.4$ & $t=3.9$\\
\midrule
$Q_G$ & total & $[-1.140,-0.960]$ & $[0.630,1.029]$\\
$\dot E$ & total & $[-1.379,-1.215]$ & $[0.406,0.771]$\\
$\dot E$ & strong cluster & $[-2.468,-2.320]$ & $[-1.273,-0.939]$\\
$\dot E$ & weak cluster & $[1.055,1.138]$ & $[1.601,1.788]$\\
\bottomrule
\end{tabular}
\end{center}
Each cluster term uses the same full-network residual and the normalization
$1/N$; this is a signed field decomposition, not a retraining ablation.
The same interval calculation certifies the strong mass and cone statements.
Finally, the speed bounds on the first and last static balls permit a
$10^{-5}$-time forward extension by a first-exit argument. Probe gates remain
constant throughout the corresponding balls and the whole sector, so the
strict endpoint-rate signs persist on the displayed short intervals.
\end{proof}

\paragraph{Scope.}
Condition~\eqref{eq:canonical-entry} specifies a parameter neighborhood of an
in-distribution trained state; it does not assume an OOD sign or a reversal.
Membership of the exact canonical initialized flow in this neighborhood at
time $3.4$ remains unproved. Thus Theorem~\ref{thm:canonical-local-continuation}
does not turn the canonical experiment into a random-initialization theorem
and does not by itself show that specialization causes reversal.

\subsection{Full-tube timing separation and minimizer localization}

The preceding theorem certifies the reversal using three snapshots. The next
result audits the entire already-certified tube and upgrades the empirical
canonical timing relation into a conditional continuous-time theorem.

\begin{theorem}[Conditional sign separation and minimizer localization]
\label{thm:canonical-full-tube-timing}
Every global compatible Clarke flow satisfying
\eqref{eq:canonical-entry} has the following properties. Put
\[
 [g_-,g_+]=[3.5738,3.6348],\qquad
 [e_-,e_+]=[3.6384,3.7152].
\]
For every compatible training-selector realization,
\begin{align}
 Q_G&<0&&\text{a.e. on }[3.4,g_-],&
 Q_G&>0&&\text{a.e. on }[g_+,3.9],\label{eq:canonical-QG-bracket}\\
 \dot E(\xi_\alpha,\cdot)&<0&&\text{a.e. on }[3.4,e_-],&
 \dot E(\xi_\alpha,\cdot)&>0&&\text{a.e. on }[e_+,3.9]
 \quad(\alpha\in\mathcal I_{\rm can}).
 \label{eq:canonical-E-bracket}
\end{align}
Consequently every minimizer of $E(\xi_\alpha,\cdot)$ on $[3.4,3.9]$
lies in $[e_-,e_+]$, uniformly for $\alpha\in\mathcal I_{\rm can}$.
The sign-transition and minimizer brackets are strictly separated:
\begin{equation}\label{eq:canonical-timing-gap}
 e_- - g_+=\frac{18}{5000}=0.0036>0.
\end{equation}
Throughout $[g_+,e_-]$, almost everywhere and uniformly in the sector,
\begin{equation}\label{eq:canonical-lag-margins}
 Q_G>1.21\cdot10^{-7},\qquad
 \dot E(\xi_\alpha,\cdot)<-5.44\cdot10^{-7},\qquad
 \dot E(\xi_\alpha,\cdot)-Q_G<-2.047\cdot10^{-4}.
\end{equation}
For the central probe the tighter minimizer bracket is
$[3.6394,3.7136]$, giving separation $23/5000=0.0046$.
Neither uniqueness nor continuity of a $Q_G$ zero, uniqueness of the OOD
minimizer, nor initialized reachability is asserted.
\end{theorem}

\begin{proof}
No flow calculation is rerun. On each of the $2500$ certified intervals
$[t_j,t_{j+1}]$, the nonnegative scalar comparison radius is nondecreasing,
so its terminal value bounds distance to the cubic reference throughout the
step. Let $y_j$ be the exact binary midpoint used by the timing verifier,
let $d_j$ bound the cubic reference motion about the midpoint, and let
$\epsilon_j$ enclose midpoint rounding. Then, with outward rounding,
\[
 \|\theta(t)-y_j\|\le R_j:=r_{j+1}+d_j+\epsilon_j
 \qquad(t_j\le t\le t_{j+1}).
\]
On this static ball Proposition~\ref{prop:canonical-static-neighborhood}
gives, for either potential,
\begin{equation}\label{eq:canonical-static-rate-bound}
 |g(\theta)^TF(\theta)-g(y_j)^TF(y_j)|
 \le R_j\{H_j(\|F(y_j)\|+L_jR_j)+\|g(y_j)\|L_j\}
       +(\|g(y_j)\|+H_jR_j)B_j,
\end{equation}
where $L_j,B_j,H_j$ are the corresponding stepwise instances of
$L_C,B_C,H_g$ from
\eqref{eq:canonical-LC}--\eqref{eq:canonical-BC}. Thus $L_j$ controls the
polynomial field, $B_j$ encloses all compatible training-gate corrections,
and $H_j$ bounds the derivative of the relevant potential gradient. For the discrepancy
$\dot E-Q_G$, the same inequality is applied directly to $g_E-g_G$ rather
than subtracting two independently widened intervals.

Probe signs are certified over every entire tube. For a reference
preactivation $q$ on a cubic Hermite step of length $h=1/5000$, the four
Bernstein control coefficients are
\[
 q_0,\qquad q_0+\frac h3\dot q_0,\qquad
 q_1-\frac h3\dot q_1,\qquad q_1.
\]
Their interval convex hull bounds the exact reference cubic. After subtracting
the parameter-tube perturbation and the $10^{-6}$-radian sector transport,
the minimum probe-gate margin over all steps and neurons is greater than
$1.40926\cdot10^{-5}$. Thus no probe gate is unresolved. The rate transport
across the sector is bounded by
\[
 2(1+\bar P_j^2/2)\bar P_j\sup\|F\|\,10^{-6},
\]
where $\bar P_j$ bounds the parameter norm on the step. Since $Q_G$ is
probe-independent, the same angular cost controls the discrepancy.

All $2500$ interval enclosures pass. The strict negative-prefix and
positive-suffix margins, rounded outward, are
\begin{center}
\begin{tabular}{lcc}
\toprule
Quantity & negative prefix: $-\text{upper}$ & positive suffix: lower\\
\midrule
$Q_G$ & $>1.0839\cdot10^{-7}$ & $>1.2158\cdot10^{-7}$\\
$\dot E$, central probe & $>3.2526\cdot10^{-7}$ & $>3.5073\cdot10^{-7}$\\
$\dot E$, entire sector & $>5.4476\cdot10^{-7}$ & $>1.9957\cdot10^{-8}$\\
\bottomrule
\end{tabular}
\end{center}
The unresolved sign steps lie entirely inside the stated brackets. There are
exactly $18$ complete steps between $g_+$ and $e_-$ on which both required
strict signs hold. On these steps the directly enclosed central discrepancy
lies in
$[-3.05302\cdot10^{-4},-2.08690\cdot10^{-4}]$; after angular transport its
upper bound remains below $-2.04714\cdot10^{-4}$.

Finally, $E$ is absolutely continuous and hence continuous on the compact
window. Integrating the strict negative prefix excludes a minimizer before
$e_-$, and integrating the positive suffix excludes one after $e_+$. The
$Q_G$ bounds similarly give a negative-to-positive sign-transition bracket
without invoking an intermediate-value theorem for a discontinuous training
field. This proves the theorem.
\end{proof}

\begin{remark}[Relation to the empirical timing law]
The theorem certifies a qualitative timing order at the canonical setting,
not the interpolated empirical lag itself. The empirical zeros near
$3.6018$ for $Q_G$ and $3.6726$ for the full derivative lie inside the
certified brackets, but they are not theorem constants. The nine-cell and
multiseed timing studies in Appendix~\ref{app:empirical} remain empirical.
\end{remark}

\subsection{Exact fixed-pattern moment dynamics}
\label{app:pattern-quotient}

The canonical reachability problem motivates a quotient that removes the
large within-pattern parameter redundancy. The following identities are exact
and independent of the computer-assisted timing proof. In this subsection use
the residual convention
\[
 \rho_n=f_\theta(x_n)-x_n
\]
and, for a training activation pattern $g\in\{0,1\}^N$, define
\[
 A_g=\frac1N\sum_{n=1}^N g_n\rho_nx_n^\top.
\]

\begin{lemma}[Fixed-pattern moment dynamics]
\label{lem:pattern-moment-dynamics}
On an open time interval on which every training preactivation is nonzero and
the neuronwise training patterns are fixed, let
\[
 I_g=\{i:\mathbf1_{u_i^\top x_n>0}=g_n\text{ for every }n\},
\]
and define
\[
 M_g=\sum_{i\in I_g}w_i u_i^\top,\qquad
 U_g=\sum_{i\in I_g}u_i u_i^\top,\qquad
 W_g=\sum_{i\in I_g}w_i w_i^\top.
\]
Then
\begin{align}
 \dot M_g&=-A_gU_g-W_gA_g,\label{eq:pattern-M}\\
 \dot U_g&=-A_g^\top M_g-M_g^\top A_g,\label{eq:pattern-U}\\
 \dot W_g&=-A_gM_g^\top-M_gA_g^\top,\label{eq:pattern-W}
\end{align}
and the training outputs satisfy
\begin{equation}\label{eq:pattern-output}
 f(x_n)=\sum_g g_nM_gx_n.
\end{equation}
Thus the group moments form an exact closed system on every fixed-pattern
cell; no concentration or low-rank approximation is used.
\end{lemma}

\begin{proof}
For $i\in I_g$, differentiation of the empirical loss gives
\[
 \dot w_i=-A_gu_i,\qquad \dot u_i=-A_g^\top w_i.
\]
Therefore
\[
 \frac{d}{dt}(w_iu_i^\top)
 =-A_gu_iu_i^\top-w_iw_i^\top A_g,
\]
and summation proves \eqref{eq:pattern-M}. Differentiating
$u_iu_i^\top$ and $w_iw_i^\top$ gives
\eqref{eq:pattern-U}--\eqref{eq:pattern-W}. Finally,
$(u_i^\top x_n)_+=g_nu_i^\top x_n$ on the cell, proving
\eqref{eq:pattern-output}. Hence every $A_g$ depends only on the data, the
fixed patterns and the group moments, establishing closure.
\end{proof}

Equivalently, set
\[
 K_g=\begin{pmatrix}U_g&M_g^\top\\M_g&W_g\end{pmatrix},\qquad
 B_g=\begin{pmatrix}0&A_g^\top\\A_g&0\end{pmatrix}.
\]
Then
\begin{equation}\label{eq:pattern-Lyapunov}
 K_g\succeq0,\qquad
 \dot K_g=-B_gK_g-K_gB_g,
\end{equation}
and $\operatorname{tr}(U_g-W_g)$ is constant on the fixed-pattern interval.
The same equations hold after refining a group by any fixed cohort label,
because every subgroup with the same training pattern uses the same $A_g$.

\begin{proposition}[Unaugmented group moments do not determine probe output or event transfer]
\label{prop:pattern-information-loss}
The collection $(M_g,U_g,W_g)$ does not in general determine an OOD probe
output, even when the individual probe signs are specified, and does not
determine the rank-one state transferred when a neuron changes training
pattern.
\end{proposition}

\begin{proof}
For the probe statement, take one training vector $e_1$, $a,b>0$, and two
balanced neurons $w_i=u_i$. Compare
\[
 u_1=(a,b),\quad u_2=(a,-b)
\]
with
\[
 u'_1=(7a/5,b/5),\quad u'_2=(a/5,-7b/5).
\]
The two pairs are related by an orthogonal mixing of neuron columns. Both
training patterns are active, and in both states the probe $-e_2$ is inactive
on neuron $1$ and active on neuron $2$. All three group moments are identical,
equal to $\operatorname{diag}(2a^2,2b^2)$, yet
\[
 f(-e_2)=(ab,-b^2),\qquad
 f'(-e_2)=(7ab/25,-49b^2/25).
\]
Thus a probe-refined moment or equivalent signed observable is necessary.

For an isolated training-gate event, put
$y_i=(u_i,w_i)$ and $K_i=y_iy_i^\top$. If neuron $i$ moves from group $g$ to
$g'$, continuity of the underlying parameters gives exactly
\[
 K_g^+=K_g^- -K_i,\qquad K_{g'}^+=K_{g'}^-+K_i.
\]
But an aggregate $K_g$ admits many rank-one decompositions, including the
orthogonal mixing above, so it does not identify $K_i$. Consequently there is
no finite universal constant $C$ for which
\[
 \|K_i-\widehat K_i\|\le C\|K_g-\widehat K_g\|
\]
holds solely from aggregate group-moment error. Event continuation therefore
requires retaining the boundary neuron's rank-one state or equivalent
geometry.
\end{proof}

\begin{remark}[Current initialized-reachability frontier]
The exact pattern-moment quotient explains why full parameter-space
coercivity is unavailable, but Proposition~\ref{prop:pattern-information-loss}
shows that the unaugmented quotient alone cannot close initialization-to-entry.
The remaining canonical obligation is to propagate enough boundary-neuron and
probe/cohort-refined information from the exact archived initialization to a
sufficient observable entry set at time $3.4$. No such initialized entry
certificate is claimed here. This limitation does not affect
Theorems~\ref{thm:canonical-local-continuation} and
\ref{thm:canonical-full-tube-timing}, which are conditional only on
\eqref{eq:canonical-entry}.
\end{remark}

\section{Additional general OOD results}
\label{app:general-theory}

The results in this section are independent supporting theory. They are not
needed for the proof of Theorem~\ref{sbsharp:theorem}, but they isolate
general features that are obscured by the explicit constants of the
finite-sample witness.

\subsection{Three snapshots certify a whole reversing sector}

\begin{proposition}[Three snapshots give a whole sector with a rebound]
\label{sbaudit:secants}
Let a finite network follow an absolutely continuous parameter trajectory
on $[t_0,t_2]$. Suppose $t_0<t_1<t_2$ and a unit $\xi_0\in\mathbb S^{d-1}$
satisfies
\begin{equation}\label{sbaudit:two-gaps}
 E(\xi_0,t_0)-E(\xi_0,t_1)\ge a>0,\qquad
 E(\xi_0,t_2)-E(\xi_0,t_1)\ge b>0.
\end{equation}
Put $L=\max_{j=0,1,2}\sum_i\|w_i(t_j)\|\|u_i(t_j)\|$.
For every unit $\xi\in\mathcal S$ with
\[
 \|\xi-\xi_0\|<\frac{\min(a,b)}{4(1+L)^2},
\]
the two gaps are at least $a/2$ and $b/2$. Consequently the error has an
interior minimum on $[t_0,t_2]$, and its rebound from that minimum to $t_2$
is at least $b/2$. The derivative is negative on a set of positive measure
in $(t_0,t_1)$ and positive on a set of positive measure in $(t_1,t_2)$.
No activation pattern need persist between snapshots.
\end{proposition}
\begin{proof}
At a snapshot the map $f$ is $L$-Lipschitz and positively homogeneous.
On unit directions its residual norm and residual Lipschitz constant are
both at most $1+L$. Thus
$|E(\xi,t_j)-E(\xi_0,t_j)|\le(1+L)^2\|\xi-\xi_0\|$.
Apply this twice to each gap. Continuity gives an attained minimum,
and both endpoints have strictly larger error than $t_1$. Absolute
continuity and integration of the derivative give the two sign sets.
The minimum need not be unique or strict; no such claim follows from
the endpoint inequalities alone.
\end{proof}

\subsection{Initial OOD descent from random initialization}
\label{app:initial-descent}


The finite-sample reversal theorem in Appendix~\ref{app:finite} uses an
explicit small-noise, large-width parameter regime. The next result isolates
a substantially more general fact: at random isotropic initialization,
strict initial OOD improvement follows with high probability whenever the
probe has a positive geometric alignment quantity \(J_\xi\). It applies to
arbitrary finite datasets and arbitrary dimension. Its sufficient finite-width
and initialization-scale conditions are asymptotic and should not be
interpreted as certifying the canonical \(h=200\) experiment.


\begin{theorem}[High-probability initial descent for a fixed OOD probe]
\label{sbaudit:initial-descent}
Fix nonzero empirical training samples in $\mathcal S$, a unit probe
$\xi\in\mathcal S$, and $d\ge2$. Initialize independently
$u_i=w_i=\varepsilon\widehat u_i$ with $\widehat u_i$ uniform on
$\mathbb S^{d-1}$. Put $s_0=h\varepsilon^2$,
$C_x=\mathbb E_n\|x\|^2$, and
\[
 \kappa_1(c)=\sqrt{1-c^2}+(\pi-\arccos c)c,\qquad
 J_\xi=\frac1\pi\mathbb E_n\left[
 (\xi^\top x)\|x\| \kappa_1\left(\frac{\xi^\top x}{\|x\|}\right)\right].
\]
For $\ell=\log(1/\delta)$ define
\[
 b_h=C_x\sqrt{\frac{6\ell}{d(d+2)h}}+
                 \frac{4C_x\ell}{3h}.
\]
With probability at least $1-\delta$ over initialization,
\begin{equation}\label{sbaudit:initial-rate}
 \dot E(\xi,0)\le-\frac{s_0}{d}
 \{J_\xi-2db_h-2dC_xs_0(2+s_0)\}.
\end{equation}
In particular, if $J_\xi\ge J_*>0$ and
\begin{equation}\label{sbaudit:initial-budget}
 2db_h+2dC_xs_0(2+s_0)\le J_*/2,
\end{equation}
then $\dot E(\xi,0)\le-s_0J_*/(2d)<0$. Almost surely all the finitely
many initial training and probe margins are nonzero, so the strict
decrease persists on some positive time interval and a nonempty open
sector containing $\xi$.
\end{theorem}
\begin{proof}
At initialization, $\|f_0(x)\|\le s_0\|x\|$ and
$\|\mathsf K_0(\xi,x)\|\le2s_0\|x\|$. Replacing both residuals by
their targets in \eqref{sbaudit:alignment} therefore costs at most
\[
 2C_xs_0^2(2+s_0).
\]
Because $w_i=u_i$, the two kernel contributions to the target pairing
are identical. With
\[
 Y_i=\mathbb E_n[(\xi^\top x)(\widehat u_i^\top\xi)_+
                                      (\widehat u_i^\top x)_+],
\]
the target pairing is $2\varepsilon^2\sum_iY_i$.
Rotational integration gives
\[
 \mathbb E[(\widehat u^\top\xi)_+(\widehat u^\top x)_+]
 =\frac{\|x\|}{2\pi d}\kappa_1(\xi^\top x/\|x\|),
 \qquad \mathbb EY_i=J_\xi/(2d).
\]
For completeness this integral follows by writing a standard Gaussian
as its independent radius and uniform direction, using
$\mathbb E\|g\|^2=d$, and integrating the product of two positive
linear forms over their angular overlap in the two-dimensional plane.
The uniform angular integral is
$\{\sin\theta+(\pi-\theta)\cos\theta\}/(4\pi)$; the independent
two-dimensional Gaussian squared radius has mean $2$.

Also $|Y_i|\le C_x$. Apply Jensen with weights proportional to
$\|x\|^2$, and use
\[
 \mathbb E[(\widehat u^\top a)^2(\widehat u^\top b)^2]
 =\frac{1+2(a^\top b)^2}{d(d+2)}\le\frac3{d(d+2)}
\]
for unit $a,b$. It follows that
$\mathbb EY_i^2\le3C_x^2/[d(d+2)]$.
The centered summands have absolute bound $2C_x$, so the one-sided
Bernstein inequality gives $h^{-1}\sum_iY_i\ge\mathbb EY_i-b_h$
with the stated probability. Substitution proves
\eqref{sbaudit:initial-rate}. On the probability-one nonzero-margin
event the finite empirical vector field is smooth near initialization;
continuity gives the last assertion.
\end{proof}

\begin{corollary}[An explicit noisy SIM region for the decreasing side]
\label{sbaudit:noisy-descent}
Consider equal cluster weights and write
$x^{(k)}=\mu_ke_k+\zeta^{(k)}$, with $\mu_1>\mu_2>0$.
Fix $0<r\le1/4$ and the OOD probe
$\xi_r=re_1-\sqrt{1-r^2}e_2$. Suppose each noise vector has norm at most
$a$, and set
\[
 \overline C=(\mu_1+a)^2,\qquad
 J_*:=\frac{11\mu_1^2r}{32\pi}-2(2\mu_1+a)a>0.
\]
Sufficient primitive initialization conditions are
\begin{equation}\label{sbaudit:primitive-descent}
 h\ge\max\left\{
 \frac{1536\overline C^2\ell}{J_*^2},
 \frac{64d\overline C\ell}{3J_*}\right\},\qquad
 h\varepsilon^2\le\min\left\{1,\frac{J_*}{24d\overline C}\right\}.
\end{equation}
Under these conditions Theorem~\ref{sbaudit:initial-descent} applies.
For planar isotropic Gaussian cluster noise with standard deviations
at most $\sigma_{\max}$ and $N$ samples in total, one may take
$a=\sigma_{\max}\sqrt{2\log(N/\delta_D)}$ with probability at least
$1-\delta_D$. Independently sampling the initialization gives total
success probability at least $1-\delta_D-\delta$.
This is an explicit nonempty region for initial descent, not yet a common
parameter region for specialization and reversal.
\end{corollary}
\begin{proof}
Write $k=\sqrt{1-r^2}$. The noiseless quantity is
\[
 J^{(0)}_{\xi_r}=\frac1{2\pi}
 \{\mu_1^2r[k+(\pi-\arccos r)r]
       -\mu_2^2k[r-k\arcsin r]\}.
\]
Here $k\ge3/4$ and $r-k\arcsin r\le r(1-k)\le r^3$.
Since $\mu_2<\mu_1$ and $r^2\le1/16$, this gives
$J^{(0)}_{\xi_r}\ge11\mu_1^2r/(32\pi)$.

To bound the noise change without differentiating the angular formula,
use the expectation representation in the preceding proof. Cauchy--Schwarz
and $\mathbb E(\widehat u^\top v)^2=\|v\|^2/d$ give, for any $x,y$,
\[
 2d\left|\mathbb E_{\widehat u}
 \bigl[(\xi^\top x)(\widehat u^\top\xi)_+(\widehat u^\top x)_+
       -(\xi^\top y)(\widehat u^\top\xi)_+(\widehat u^\top y)_+\bigr]
 \right|
 \le2(\|x\|+\|y\|)\|x-y\|.
\]
Average this inequality with $y=\mu_ke_k$ to get $J_{\xi_r}\ge J_*$.
The two width bounds make the two terms of $2db_h$ at most $J_*/8$
each. The initialization cap makes the remaining term in
\eqref{sbaudit:initial-budget} at most $J_*/4$. The Gaussian assertion
is the two-dimensional radial tail followed by a union bound. Finally,
choose sufficiently small positive noise, then finite $h$, then positive
$\varepsilon$ sufficiently small to see nonemptiness for this corollary.
\end{proof}

\begin{corollary}[A directional noise estimate reaches noise-scale probe arcs]
\label{sbaudit:sharp-noise}
Let $0<r_0<r_1\le1/4$, suppose all sample perturbations have norm at most
$a\le\mu_2/4$, and suppose the empirical mean perturbation of cluster 1
has norm at most $b$. Uniformly for $r_0\le r\le r_1$,
\begin{equation}\label{sbaudit:sharp-J}
 J_{\xi_r}\ge J_{\rm sh}:=\frac1{2\pi}
 \left\{\mu_1(\mu_1-a)r_0-\mu_1b-a^2
              -\frac{(\mu_2r_1+a)^3}{\mu_2-a}\right\}.
\end{equation}
Whenever $J_{\rm sh}>0$, it may replace $J_*$ in the initialization
conditions, with $\overline C=(\mu_1+a)^2$.
For $n_1$ independent isotropic Gaussian samples in cluster 1, the mean
event can be imposed at failure probability $\delta_m$ with
\[
 b=\sigma_1\sqrt{2\log(1/\delta_m)/n_1}.
\]
\end{corollary}
\begin{proof}
The function $\kappa_1$ in Theorem~\ref{sbaudit:initial-descent} is increasing,
with $\kappa_1(0)=1$, so $c\kappa_1(c)\ge c$ for every
$-1\le c\le1$.
On cluster 1 this gives a lower bound by
$\mathbb E_{1,n}[(\xi_r^\top x)\|x\|]$ for the numerator of its
contribution to $J_{\xi_r}$. Since
$|\|x\|-\mu_1|\le a$ and $|\xi_r^\top x|\le\mu_1r+a$,
this is at least
$\mu_1^2r-\mu_1b-\mu_1ra-a^2$.

On cluster 2, $\xi_r^\top x<0$. Write its angle with $\xi_r$ as
$\pi-\beta$, with $0\le\beta<\pi/2$. Then
\[
 \kappa_1(-\cos\beta)=\sin\beta-\beta\cos\beta
 \le\sin\beta(1-\cos\beta)\le\sin^3\beta.
\]
Moreover $|\det(\xi_r,x)|\le\mu_2r+a$ and $\|x\|\ge\mu_2-a$.
Its possibly negative numerator is therefore bounded below by
$-(\mu_2r+a)^3/(\mu_2-a)$. Combine the two clusters with weights
$1/2$, and use $r\ge r_0$ in the positive term and $r\le r_1$ in the
negative term. The mean-event assertion follows from the planar
Gaussian radial tail with variance $\sigma_1^2/n_1$.
\end{proof}

\begin{corollary}[Uniform descent permits a later, trajectory-chosen probe]
\label{sbaudit:uniform-descent}
Fix $0<r_0<r_1\le1/4$ and the closed probe arc
$\mathcal I=\{re_1-\sqrt{1-r^2}e_2:r_0\le r\le r_1\}$.
Use a positive uniform lower bound $J_*$ from Corollary~\ref{sbaudit:noisy-descent}
(evaluated at $r_0$) or from \eqref{sbaudit:sharp-J}.
Let
\[
 \rho=\frac{J_*}{32d\overline C},\quad
 G=1+\left\lceil\frac{\arcsin r_1-\arcsin r_0}{\rho}\right\rceil,
 \qquad \ell=\log(G/\delta).
\]
Impose \eqref{sbaudit:primitive-descent} with this $\ell$.
With initialization probability at least $1-\delta$, every probe on
$\mathcal I$ has initial upper right error derivative at most
$-s_0J_*/(4d)$. There exists a common $\tau>0$ on which all these errors
are strictly decreasing. Thus a later probe chosen using the trajectory,
anywhere in this deterministic arc, is covered by the same event.
The positive time $\tau$ is not quantitatively lower bounded here.
\end{corollary}
\begin{proof}
Use $G$ equally spaced angular grid points. Every point on the arc is
within Euclidean distance $\rho$ of a grid point. The Bernstein argument
holds at every grid point simultaneously with probability at least
$1-\delta$. Directly from its definition,
\[
 |Y_i(\xi)-Y_i(\zeta)|\le2C_x\|\xi-\zeta\|
 \quad(\|\xi\|=\|\zeta\|=1).
\]
Consequently the transport cost in $2d h^{-1}\sum_iY_i$ is at most
$4dC_x\rho\le J_*/8$. Grid fluctuations cost at most $J_*/4$ and
the residual correction costs at most $J_*/4$. The remaining margin
is at least $3J_*/8$, stronger than asserted.

At an initial zero probe preactivation the target-pairing formula is
unchanged for every compatible selector: its potentially ambiguous
second kernel term has a factor $w_i^\top\xi=u_i^\top\xi=0$.
The residual correction is uniformly bounded for all those selectors.
Training margins are nonzero almost surely, so the parameter vector
field is continuous for a positive initial interval. The graph of
compatible probe selectors is closed. Compactness of the probe arc
and of the selector cube, together with the strict uniform initial
margin, therefore gives a common positive interval on which all
compatible error slopes remain negative. Integrate the almost-everywhere
slopes to obtain strict decrease. This step provides existence of an
interval, not an explicit initialized learning clock.
\end{proof}

\begin{remark}[Canonical alignment quantity]
For the canonical SIM dataset and the \(-85^\circ\) reversing probe, the
empirically evaluated alignment quantity is \(J_\xi\approx0.148>0\).
Thus the geometric sign condition of Theorem~\ref{sbaudit:initial-descent}
holds at the probe where the reversal phenomenon is observed. The theorem's
sufficient width and initialization-scale inequalities are not claimed to
certify the experimental width \(h=200\); the comparison is only at the level
of the alignment mechanism.
\end{remark}



\subsection{An exact signed transverse-energy identity}

The following identity is retained as an additional structural result.  It is
not used to explain the canonical reversal timing, which is described by the
entrywise normal form in Appendix~\ref{app:normal}; in particular, we do not
claim that a global transverse-energy ratio contracts over the empirical
degradation window.

Let \(\mathcal S=\operatorname{span}\{e_1,e_2\}\) be a two-dimensional
concept subspace and write
\[
 v_i=P_{\mathcal S}u_i,\qquad z_i=P_{\mathcal S}w_i,\qquad
 f_\parallel=P_{\mathcal S}f,\qquad e=x-f_\parallel.
\]
For compatible training selectors \(\alpha_i\), define
\[
 A_i=\mathbb E_n[e x^\top\alpha_i],\qquad
 \dot z_i=A_iv_i,\qquad
 \dot v_i=A_i^\top z_i-P_i,
\]
where
\[
 P_i=\mathbb E_n[x\alpha_i\langle w_{i,\perp},f_\perp\rangle].
\]
Fix a partition \([h]=J_1\sqcup J_2\), let \(p(i)\) be the assigned concept
axis and \(q(i)=3-p(i)\), and set
\[
 T=\sum_i(v_{i,q(i)}^2+z_{i,q(i)}^2),\qquad
 Y=\sqrt T,\qquad
 \Pi=\Big(\sum_i\|P_i\|^2\Big)^{1/2}.
\]

\begin{lemma}[Signed transverse-energy identity]
\label{sbsign:energy}
Define
\[
 f_D(x)=\sum_i\alpha_i(x)
   \bigl(z_{i1}v_{i1}x_1e_1+z_{i2}v_{i2}x_2e_2\bigr),
\]
\[
 f_O(x)=\sum_i\alpha_i(x)
   \bigl(z_{i1}v_{i2}x_2e_1+z_{i2}v_{i1}x_1e_2\bigr),
\]
and
\[
 f_{TT}(x)=\sum_i\alpha_i(x)
 z_{i,q(i)}v_{i,q(i)}x_{q(i)}e_{q(i)}.
\]
Then, almost everywhere,
\begin{equation}\label{sbsign:identity}
 \frac12\dot T
 =-\|f_O\|_n^2+\langle x-f_D,f_O\rangle_n
 +2\sum_i A_{i,q(i)q(i)}z_{i,q(i)}v_{i,q(i)}
 -\sum_i v_{i,q(i)}P_{i,q(i)}.
\end{equation}
Writing \(m_\times=\mathbb E_n|x_1x_2|\), and, for \(T>0\),
\[
 \kappa(t)=
 \frac{2\sum_i A_{i,q(i)q(i)}z_{i,q(i)}v_{i,q(i)}
       -\|f_O\|_n^2}{T},
\]
with \(\kappa=0\) when \(T=0\), one has
\begin{equation}\label{sbsign:Y}
 \dot Y\le
 \kappa Y+\sqrt{2S}(1+S)m_\times+\Pi.
\end{equation}
If \(\|A_i\|\le K\), then
\[
 |\kappa|\le K+2\lambda_nS,
 \qquad
 \lambda_n:=\max\{\mathbb E_nx_1^2,\mathbb E_nx_2^2\}.
\]
\end{lemma}

\begin{proof}
Differentiate \(T/2\) using the projected equations.  The output-coordinate
and input-coordinate terms together give
\(\langle e,f_O+2f_{TT}\rangle_n\), while the ambient term contributes
\(-\sum_i v_{i,q(i)}P_{i,q(i)}\).  Substituting
\(e=x-f_D-f_O\) in the first inner product gives
\eqref{sbsign:identity}.

Write
\(f_D=(d_1(x)x_1,d_2(x)x_2)\) and
\(f_O=(o_{12}(x)x_2,o_{21}(x)x_1)\).
Balance and Cauchy--Schwarz give \(|d_r(x)|\le S\).  If
\[
 Q_{12}=\sum_i|z_{i1}v_{i2}|,\qquad
 Q_{21}=\sum_i|z_{i2}v_{i1}|,
\]
then
\[
 Q_{12}+Q_{21}
 \le
 \Big(\sum_i(v_{i,p(i)}^2+z_{i,p(i)}^2)\Big)^{1/2}\sqrt T
 \le\sqrt{2ST}.
\]
Hence
\[
 |\langle x-f_D,f_O\rangle_n|
 \le(1+S)m_\times\sqrt{2ST},
\]
and the ambient term is at most \(\sqrt T\,\Pi\).
Dividing by \(Y\) proves \eqref{sbsign:Y}, with the usual regularization
at \(Y=0\).
Finally,
\[
 \|f_O\|_n^2
 \le \lambda_n(Q_{12}^2+Q_{21}^2)
 \le2\lambda_nST,
\]
which yields the stated bound on \(\kappa\).
\end{proof}

\begin{remark}[Scope of the identity]
Lemma~\ref{sbsign:energy} isolates an exact negative contribution
\(-\|f_O\|_n^2\), but the remaining forcing terms need not be small at the
canonical noise level.  We therefore use it only as a structural identity,
not as the explanation of the empirical reversal.  The canonical mechanism
is instead summarized by the specific block entries appearing in
\(\mathcal G\) in Appendix~\ref{app:normal}.
\end{remark}


\section{Finite-sample OOD non-monotonicity}\label{app:finite}

\subsection{The finite-sample theorem}

\begin{theorem}[Finite-sample coherent learning with OOD non-monotonicity]
\label{sbsharp:theorem}
Let $d=2$, $h\ge10^7$, $\varepsilon=\sqrt{0.04/h}$, and let $N=2000$
consist of independent samples with exactly half in each isotropic
Gaussian cluster, with
\[
 \mu_1=1,\qquad\mu_2=10^{-4},\qquad\sigma_1=\sigma_2=10^{-9}.
\]
Use independent isotropic balanced initialization. With probability at
least $0.97$, one fixed cohort has at least $0.24h$ neurons, input and
output angles at most $0.13$ radians, and longitudinal learned mass at
least $0.55$ throughout $[t_{\rm spec},t_b]$. With reference levels
$A^2=0.69,0.71,0.715,0.95$ defining $t_-,t_{\rm spec},t_a,t_b$,
\[
 0<0.00017<0.0055<t_-<t_{\rm spec}<t_a<t_b<7.6.
\]
For every $0\le t\le t_b$, $S(t)<2$ and $|f_2(e_2,t)|<0.011$.
On the entire deterministic open arc
\[
 \xi_\theta=(\sin\theta,\cos\theta),\qquad
 3.5\cdot10^{-5}<\theta<4\cdot10^{-4},
\]
whose angular width is $3.65\cdot10^{-4}$ radians,
\begin{align*}
 \dot E(\xi_\theta,t)&\le-1.5\cdot10^{-9}
                         &&(0\le t\le0.00017),\\
 E(\xi_\theta,0)-E(\xi_\theta,0.00017)&>2.04\cdot10^{-12},\\
 E(\xi_\theta,t_b)-E(\xi_\theta,t_a)&>7.2\cdot10^{-6},\\
 \dot E(\xi_\theta,t)&>3.56\cdot10^{-6}
                         &&(t_-\le t\le t_b)
\end{align*}
almost everywhere where derivatives are stated. On the inner arc
$3.5\cdot10^{-5}<\theta<4.5\cdot10^{-5}$ the rebound exceeds
$2\cdot10^{-5}$. On the upper subarc
$3\cdot10^{-4}<\theta<4\cdot10^{-4}$, of width $10^{-4}$,
\[
 \dot E\le-2.67\cdot10^{-8}\quad(0\le t\le0.0055),\qquad
 E(\xi,0)-E(\xi,0.0055)>2.15\cdot10^{-9}.
\]
It retains the full-arc rebound $7.2\cdot10^{-6}$.
Every global minimum
on $[0,t_b]$ belongs to $(0,t_-]$, hence occurs strictly before the
certified specialization time. The statement certifies an initial
transient and subsequent degradation of a learned dominant-family state;
it does not locate a reversal caused by specialization.
\end{theorem}

\begin{remark}[Scope of the certified regime]
The theorem is an existence-and-robustness result for OOD non-monotonicity,
not a quantitative theorem for the canonical experiment of
Appendix~\ref{app:empirical}. Its explicit witness is deliberately
small-noise and highly asymmetric. A visible proof bottleneck is the
weak-cluster perturbation budget
\[
 P=30\{a_D+(\mu_2+a_D)^2\},
\]
together with concentration and transport charges that must be small
relative to the weak scale $\mu_2$. These requirements force a regime far
from the canonical values before the later comparison argument even begins.
The large width is likewise a conservative finite-width concentration
requirement. Closing this gap while deriving the empirical
$\mathcal G$-crossing relation from random initialization remains open.
\end{remark}

The proof combines a single concentration event, a relative comparison with
logistic growth, and signed bounds for the full probe output.
Appendix~\ref{app:normal} is not a premise of this theorem.

\subsection{A common high-probability event}
Write $\widehat u_i=(t_i,y_i)$, $Q=\sum_i(t_i)_+^2$,
$A_0=\varepsilon\sqrt Q$, $v_i=(t_i)_+/\sqrt Q$, and
$I_C=\{i:t_i\ge|y_i|\}$. Define
\[
 K_0=\varepsilon^2\sum_i(y_i)_+^2,\quad
 C_0=\varepsilon^2\sum_i t_i(y_i)_+,\quad b_0=|B(0)v|,
 \quad B(0)=(\varepsilon y_i)_i.
\]
\begin{proposition}[Initialization and sample event]\label{prop:good-event}
At the parameters of Theorem~\ref{sbsharp:theorem}, with probability greater
than $1-0.020575$, all of the following hold:
\begin{gather}\label{eq:good-moments}
 Q/h,\ K_0/s\in[0.248,0.252],\qquad
 \left|h^{-1}\sum_i(t_i)_+(y_i)_+-\frac1{4\pi}\right|\le0.002,\\
 |C_0|\le0.001s,\quad |I_C|\ge0.24h,\quad
 h^{-1}\sum_{i\in I_C}t_i^2\ge\frac18+\frac1{4\pi}-0.002,
 \quad b_0\le1.25\varepsilon.\nonumber
\end{gather}
Every sample is within $a_D=10^{-8}$ of its cluster center.
At $e_1,e_2$ and the eleven equally spaced angles in
$[3.5\cdot10^{-5},4.5\cdot10^{-5}]$, the initial output satisfies
\begin{equation}\label{eq:good-net}
 \|f_0(v)-(s/4)v\|\le d_{\rm net}:=
 \sqrt2s\left(\sqrt{\frac{\ell}{4h}}+\frac{4\ell}{3h}\right),
 \qquad \ell=\log10400.
\end{equation}
The residual map $v\mapsto f_0(v)-(s/4)v$ is $0.05$-Lipschitz on
unit directions. In particular its error on normalized training samples
is at most $d_{\rm train}=d_{\rm net}+0.05(2a_D/\mu_2)$.
\end{proposition}
\begin{proof}
Uniform-circle integration gives means $1/4,1/4,1/(4\pi),0,1/4$
and $1/8+1/(4\pi)$ for the six averages in
\eqref{eq:good-moments}, in their stated order. Their ranges have lengths
$1,1,1/2,1,1,1$. Hoeffding's inequality and a union bound give failure at most
\[
 6e^{-2h(0.002)^2}+2e^{-8h(0.002)^2}
 +2e^{-2h(0.001)^2}+2e^{-2h(0.01)^2}<4.13\cdot10^{-9}.
\]
For $b_0$, use the independent variables
$X_i=t_i y_i\mathbf1_{\{t_i>0\}}$, not independence of $B(0)$ and $v$.
They have mean zero, variance $1/16$ and third absolute moment $1/(12\pi)$.
The iid Berry--Esseen inequality with constant $1/2$ (a conservative
consequence of \cite{shevtsova2011}) therefore bounds the two
CDF errors by $16/(3\pi\sqrt h)\le8\sqrt2/(\pi\sqrt h)$.
On $Q\ge0.248h$, the event $b_0>1.25\varepsilon$ requires
$|4\sum_iX_i/\sqrt h|>x_0:=5\sqrt{0.248}$.
The Gaussian Mills bound gives the additional failure bound
\[
 \frac{2e^{-x_0^2/2}}{x_0\sqrt{2\pi}}+
 \frac{8\sqrt2}{\pi\sqrt h}<0.015575.
\]
For fixed unit $v$, the scalar
$Y_j=\widehat u_j(\widehat u^\top v)_+$ has mean $v_j/4$,
variance $(1+v_j^2)/16\le1/8$, and centered absolute value at most two.
Two-sided Bernstein at thirteen directions and two coordinates costs
$52e^{-\ell}=0.005$, with Euclidean error \eqref{eq:good-net}.
ReLU Lipschitzness bounds the residual map's Lipschitz constant by
$s+s/4=0.05$. The distance of a normalized sample from its center
direction is at most $2a_D/\mu_p$, proving the training bound.
Finally, for planar isotropic noise,
$\Pr(\|\xi\|>a_D)=\exp[-a_D^2/(2\sigma^2)]=e^{-50}$;
a union bound costs $2000e^{-50}$. Summing these failures is less than
$0.020575$. The events need not be independent.
\end{proof}

\begin{remark}[Why the theorem states $0.97$]
The explicit arithmetic ledger evaluates the total failure budget as
$0.0205743073$, so the displayed union bound actually gives success
probability above $0.9794$; the theorem rounds this down to $0.97$.
The $b_0$ tail contributes $0.0155743032$, about $75.7\%$ of the total
failure budget, so the probability constant is dominated by this
conservative one-dimensional tail estimate rather than by the dynamical
comparison. In the seed-$20260918$ numerical illustration,
$b_0=8.12\cdot10^{-7}$ while the certified ceiling is
$7.91\cdot10^{-5}$, a factor of about $97$. This realized numerical slack
is illustrative only and does not improve the theorem's probability
guarantee.
\end{remark}

\subsection{Logistic reference and time convention}
All comparisons in the next two subsections use $\tau=t/2$.
Set $s=0.04$, $H=3.8$, $\nu=0.05$ and
$A_{0,\pm}=\sqrt{s(0.25\pm0.002)}$.
Let $A'=A(1-A^2)$, $A(0)=A_0$. Its solution and physical hitting time are
\begin{equation}\label{eq:logistic-time}
 A(\tau)^2=\frac{A_0^2e^{2\tau}}{1-A_0^2+A_0^2e^{2\tau}},
 \qquad t(y)=\log\frac{y(1-A_0^2)}{(1-y)A_0^2}.
\end{equation}
Thus $t_-=t(0.69)$, $t_{\rm spec}=t(0.71)$, $t_a=t(0.715)$,
$t_b=t(0.95)$ are well defined and strictly ordered.
The terminal normalized time is at most
$\frac12\log[0.95(1-A_{0,-}^2)/(0.05A_{0,-}^2)]<3.773836<H$;
$t_->\log[0.69(1-A_{0,+}^2)/(0.31A_{0,+}^2)]>5$.


\subsection{Noisy dominant-axis dynamics}

Write $a,c$ for the full first-coordinate input and output vectors and
$B,D$ for the second-coordinate rows. Set
\begin{gather*}
 a_+=(a_i)_+,\qquad u=\|a_+\|,\qquad p=c-a,\qquad
 m=c^\top a_+,\qquad k=1-m,\\
 R=Da_+,\qquad z=|R|/u,\qquad l=\|p\|/u,\qquad W=z+l.
\end{gather*}
Thus the actual output at $e_1$ is $(m,R)$, without any block reduction.



\begin{lemma}[Perturbed dominant-axis dynamics]
\label{sbwide:axis-system}
Suppose every strong sample is within $a_D\le10^{-3}$ of $e_1$, every
weak sample within $a_D$ of $\mu_2e_2$, and $0<\mu_2\le10^{-2}$.
On a stopped interval with $S\le2$, $0\le k\le1$, and $|R|\le1$, put
\[
 P=30\{a_D+(\mu_2+a_D)^2\}.
\]
There is a diagonal matrix $\mathcal A\in[0,1]^{h\times h}$ such that
almost everywhere
\begin{align}\label{sbwide:perturbed-axis}
 a'&=ka_++k\mathcal A p-\mathcal A D^\top R+E_a,
 &c'&=ka_++E_c,\\
 D'&=-Ra_+^\top+E_D,
 &B'&=E_B,\nonumber
\end{align}
where every error has Frobenius norm at most $P$ and every neuronwise
error has norm at most $P\|u_i\|$. Furthermore
\[
 \|(\mathcal A-I)a_+\|\le a_D\sqrt S.
\]
No common strong-sample mask is assumed. The weak selectors are arbitrary.
\end{lemma}
\begin{proof}
Use the selected equations of Lemma~\ref{lem:gf-regularity}.
Let $\mathcal A_{ii}$ be the empirical strong-cluster mean of its actual
selector. Write $f(e_1)=(m,R)$. ReLU Lipschitzness gives
\[
 \|f(x)-f(e_1)\|\le Sa_D,\qquad
 |(u_i^\top x)_+-(a_i)_+|\le a_D\|u_i\|.
\]
The target-minus-output residual changes by at most $(1+S)a_D$.
At $e_1$ its norm is at most $\sqrt2$ on the stopped interval.
The output-gradient error on the strong cluster is therefore at most
$[3(1+a_D)+\sqrt2]a_D\|u_i\|<5a_D\|u_i\|$.
The first input-gradient error before replacing $\mathcal A a$ by $a_+$
is also less than $5a_D\|u_i\|$. A changed gate implies $|a_i|\le a_D\|u_i\|$;
thus $|\mathcal A_{ii}a_i-(a_i)_+|\le a_D\|u_i\|$.
The transverse input-gradient norm is at most
$(\sqrt2+3a_D)a_D\|u_i\|$.
On a weak sample, $\|x-f(x)\|\le3\|x\|$, so each input or output
gradient norm is at most $3(\mu_2+a_D)^2\|u_i\|$.
Corollary~\ref{cor:exact-balance} gives $\|w_i\|=\|u_i\|$ in both gradient estimates.
Combining these estimates gives each error at most
$9\{a_D+(\mu_2+a_D)^2\}\|u_i\|$.
Summing squares and using $\sqrt S\le\sqrt2$ proves both asserted
bounds with the stated reserve. The selector estimate follows from the
same changed-gate implication.
\end{proof}

Define $A,v$ by \eqref{eq:logistic-time} and Proposition~\ref{prop:good-event}.
For the next lemmas, $0<A_{0,-}\le A_0$, $0<\nu<1$, and
$A^2\le1-\nu$ on $[0,T]$, with $T\le H$.

\subsection{Weighted tracking estimates}

Set $q_c=10^{-3}$, $\beta_B=\sqrt{s}+PH$, and suppose $A^2\le1-\nu$.
Define
\begin{align*}
 I_*&=\operatorname{arctanh}\sqrt{1-\nu}
                 -\operatorname{arctanh}A_{0,-},\\
 J_*&=\tfrac12\log\frac{1-A_{0,-}^2}{\nu},\\
 Z_*&=\{b_0+PH(1+6\beta_B/A_{0,-})\}
                  e^{(\beta_B^2+0.01)H},\\
 l_*&=\beta_B Z_*e^{2q_c}\frac{I_*}{\sqrt{1-\nu}}
               +\frac{2PH e^{q_c}}{A_{0,-}},\qquad W_*=Z_*+l_*,\\
 q_*&=H\max\{l_*+l_*^2+Z_*^2+2P/A_{0,-},\,
                     l_*+\beta_BZ_*+2P/A_{0,-}\},\\
 C_c^*&=3q_*+(H+1)l_*+PH,\qquad
 D_*^{\rm mov}=Z_*e^{2q_*}J_*+PH.
\end{align*}



\begin{lemma}[Weighted residual and mismatch bounds]
\label{sbsharp:weighted}\label{sbwide:relative}
Suppose $0<\nu<1$, $\beta_B\le0.3$, and
\[
 W_*<\min\{0.005,\nu/16\},\qquad
 q_*<\min\{q_c,\nu/8\},\qquad
 \frac{3P}{A_{0,-}}<\min\{0.001,\nu/16\}.
\]
Assume also
\[
 1-\nu/2+\{\sqrt s+H[(1+\beta_B)W_*+P]\}^2+\beta_B^2<2.
\]
Then throughout $[0,T]$,
\begin{align}\label{sbwide:relative-bounds}
 z&\le Z_*,& l&\le l_*,& |\log(u/A)|&\le q_*,\\
 \|a_+/u-v\|&\le q_*,&\|a_+-Av\|&\le3Aq_*,&
 \|c-(Av+a(0)_-)\|&\le C_c^*,\nonumber\\
 \|D-B(0)\|&\le D_*^{\rm mov},&\|B-B(0)\|&\le PH,&
 \|D\|&\le\beta_B,\quad S<2.\nonumber
\end{align}
Here $a(0)_-=(\min\{a_i(0),0\})_i$.
Every neuron is included.
\end{lemma}
\begin{proof}
Use Lemma~\ref{sbwide:axis-system} and bootstrap on
$S<2$, $0<k<1$, $|R|<1$, $A_{0,-}/2<u$, $u^2<1-\nu/2$,
$z+l<0.005$ and $|\log(u/A)|<q_c$. These bounds hold initially.
Let $\mathsf D_+$ be the diagonal derivative selector of $a_+$; thus
$\mathsf D_+ a_+=a_+$ almost everywhere by
Lemma~\ref{lem:gf-regularity}. Subtraction and the product rule give exactly
\begin{align*}
 p'&=-k\mathcal A p+\mathcal A D^\top R+E_c-E_a,\\
 R'&=(k-u^2)R+kD\mathsf D_+\mathcal A p
       -D\mathsf D_+\mathcal A D^\top R+E_Da_++D\mathsf D_+ E_a.
\end{align*}
The quadratic term is dissipative, since the diagonal matrices commute
and have nonnegative entries. The equation for $D$ also gives
$\|D\|\le\sqrt s+PH=\beta_B$, and $\|B-B(0)\|\le PH$.
Put $\theta=a_+^\top p/u^2$ and $g=u'/u$.
The selector error in Lemma~\ref{sbwide:axis-system} gives
\begin{align*}
 g&=k(1+\theta)-z^2+E_g,
 &|E_g|&\le P(1+l+z)/u,\quad |\theta|\le l,\\
 g&=1-u^2+(1-2u^2)\theta-u^2\theta^2-z^2+E_g.
\end{align*}
For example the radial selector correction is bounded by
$a_D\sqrt S(k l+\beta_Bz)/u$ in addition to $P/u$;
$a_D\sqrt2\le P$ and $\beta_B<1$ give the stated bound.
The exact residual and mismatch equations imply, almost everywhere,
\begin{align*}
 l'&\le-gl+\beta_Bz+2P/u,\\
 z'&\le(-u^2-k\theta+z^2+|E_g|)z
                    +k\beta_Bl+P(1+\beta_B/u),
 \qquad |\theta|\le l.
\end{align*}
For $Z=z+\beta_Bl$ the $k\beta_Bl$ term cancels the matching part of
$-\beta_Bgl$ because $k-g=-k\theta+z^2-E_g$. Thus
\[
 Z'\le(\beta_B^2+l+z^2+|E_g|)Z+P(1+3\beta_B/u)
 \le(\beta_B^2+0.01)Z+P(1+6\beta_B/A_{0,-}).
\]
The last inequality uses the stopping bounds and
$|E_g|\le P(1+l+z)/u$. This proves $Z\le Z_*$.
Independently, dissipativity in the Cartesian mismatch equation gives
$(\|p\|)'\le\beta_Buz+2P$. Since
$\int_0^\tau A\,ds=\operatorname{arctanh}A(\tau)
-\operatorname{arctanh}A_0$, division by $u\ge Ae^{-q_c}$ gives $l_*$.
The function $(\operatorname{arctanh}A-\operatorname{arctanh}A_0)/A$
is increasing in $A\ge A_0$, justifying the uniform terminal bound.

In the exact radial identity, the forcing after $1-u^2$ is at most
$l_*+l_*^2+Z_*^2+2P/A_{0,-}$. The directional speed is at most
$l_*+\beta_BZ_*+2P/A_{0,-}$. For $q=\log(u/A)$ one has
$q'=A^2(1-e^{2q})+\epsilon_q$, with
$|\epsilon_q|\le l_*+l_*^2+Z_*^2+2P/A_{0,-}$.
The first term has sign opposite to $q$, so integration of the upper
Dini derivative of $|q|$ gives $|q|\le q_*$.
The projected equation for $a_+/u$ has speed at most
$l_*+\beta_B Z_*+2P/A_{0,-}$, giving its distance bound.
Then $\|a_+-Av\|\le A\{e^{q_*}q_*+e^{q_*}-1\}\le3Aq_*$,
because $q_*<10^{-3}$. 
These strict bounds close the $z+l$ and logarithmic faces.
For completeness, the radius bounds are
$u>A_{0,-}/2$ and $u^2<(1-\nu)e^{\nu/4}<1-\nu/2$.
Then $0<m=u^2(1+\theta)<1$, $k\ge\nu/2-l_*>0$,
$|R|\le uZ_*<1$, and
$g\ge\nu/2-(l_*+l_*^2+Z_*^2+2P/A_{0,-})>0$.

The strict smallness assumptions imply the last positivity, since
$l_*+l_*^2+Z_*^2\le W_*+W_*^2<1.005\nu/16$
and $2P/A_{0,-}<\nu/24$.

On $a_i<0$ the $ka_+$ term vanishes, so integration of
$|a_i'|\le k|p_i|+|D_iR|+|E_{a,i}|$ gives
for the negative input row
$\|a_-\|\le\sqrt s+H[(1+\beta_B)W_*+P]$.
Together with $\|B\|\le\beta_B$ and the assumed mass inequality,
this excludes $S=2$. These estimates exclude every stopping face.

For the output comparison, ReLU Lipschitzness gives
$\|c_+-a_+\|\le\|p\|\le l_*$. On $c_i<0$ one has
$(a_i)_+\le|p_i|$, so $c'=ka_++E_c$ gives
$\|c_--a(0)_-\|\le H(l_*+P)$.
Combining these estimates with $\|a_+-Av\|\le3Aq_*$ proves $C_c^*$. 
Finally $\|D'\|\le u^2z+P$ and
$\int_0^\tau A^2\,ds=\frac12\log[(1-A_0^2)/(1-A(\tau)^2)]$
give $D_*^{\rm mov}$. 
\end{proof}

\subsection{Persistent cohort and output envelopes}

Use the fixed cohort $I_C$ of Proposition~\ref{prop:good-event}.

Define the cohort charges using Lemma~\ref{sbsharp:weighted}:
\begin{align*}
 M_*&=4q_*+2W_*,& V_*&=H(M_*+W_*+P),\\
 J_D&=e^{V_*}H(W_*+P)/A_{0,-},&
 J_B&=e^{V_*}HP/A_{0,-},\\
 J_p&=H(1+J_D)W_*+2e^{V_*}HP/A_{0,-},\\
 \Gamma_C&=e^{M_*H}-1+
 \sqrt2e^{2M_*H}H\{J_p+(1+J_D)W_*+Pe^{V_*}/A_{0,-}\}.
\end{align*}



\begin{lemma}[Persistence of a coherent cohort]
\label{sbwide:core-holding}
Suppose Lemma~\ref{sbsharp:weighted} applies, $\Gamma_C<1/4$, and
\[
 \frac{1-\Gamma_C}{\sqrt2}>a_De^{V_*}.
\]
Then all strong noisy training gates of $I_C$ stay active and
\[
 a_i\ge(1-\Gamma_C)A\varepsilon t_i/A_0,\qquad
 |p_i|\le\varepsilon J_p,\qquad
 |B_i|\le\varepsilon(|y_i|+J_B),\qquad
 |D_i|\le\varepsilon(|y_i|+J_D).
\]
If $\sum_{I_C}v_i^2\ge\omega_C$, its longitudinal learned mass is at least
\begin{equation}\label{sbwide:core-mass}
 A^2\omega_C(1-\Gamma_C)^2-\sqrt{s}J_p.
\end{equation}
For any $A_{\rm spec}>0$, on $A\ge A_{\rm spec}$, its input and output tangent-angle bounds are
respectively
\[
 \frac{A_0}{A_{\rm spec}}
 \frac{1+\sqrt2J_B}{1-\Gamma_C},\qquad
 \frac{A_0}{A_{\rm spec}}
 \frac{1+\sqrt2J_D}
 {1-\Gamma_C-\sqrt2 A_0J_p/A_{\rm spec}},
\]
provided the last denominator is positive.
\end{lemma}
\begin{proof}
The neuronwise residual matrix has norm at most
$k+|R|+P$ on the stopped interval: the strong residual contributes
$k+|R|$, and the data estimates in
Lemma~\ref{sbwide:axis-system} absorb the remaining terms.
Full balance thus gives
\[
 \|u_i\|=\|w_i\|\le\varepsilon\frac A{A_0}e^{V_*}.
\]
Here $|m-A^2|\le M_*$ follows from Lemma~\ref{sbsharp:weighted},
and $\int(1-A^2)=\log(A/A_0)$.
Integrating the per-neuron $B,D,p$ equations gives $J_B,J_D,J_p$.
Stop the core only when a strong noisy gate first ceases to be active.
Before that time,
$a_i'=ka_i+kp_i-D_iR+E_{a,i}$.
The integrating factor for $k$ is between
$(A/A_0)e^{-M_*H}$ and $(A/A_0)e^{M_*H}$.
Variation of constants and $t_i\ge1/\sqrt2$ give the displayed
relative error $\Gamma_C$.
The gate reserve follows because its preactivation is at least
$a_i-a_D\|u_i\|$; the stated inequality makes it strictly positive.
For the mass bound use $c_i=a_i+p_i$ and
$\sum_{I_C}|a_i|\le\sqrt h\|a_+\|\le\sqrt h$.
Finally $|y_i|\le t_i$ on this cohort proves the two angle bounds.
\end{proof}


\begin{lemma}[Output bounds on the reference interval]\label{lem:output-envelope}
On the event of Proposition~\ref{prop:good-event}, under
Lemma~\ref{sbsharp:weighted}, define
Put $K_+=0.252s$, $C_+=0.001s$ and
\[
 P_-=\sqrt{s}\frac{1/(4\pi)-0.002}{\sqrt{0.252}},\qquad
 P_+=\sqrt{s}\frac{1/(4\pi)+0.002}{\sqrt{0.248}}.
\]
The two coordinate errors are
\[
 e_1=C_c^*\sqrt{K_+}+\sqrt2PH<0.000213365,\qquad
 e_2=D_*^{\rm mov}\sqrt{K_+}+\beta_BPH<0.000020428.
\]

For every $0\le t\le t_b$,
\[
 |f_1(e_2,t)-\{P_\xi(A-A_0)+C_0\}|\le e_1,\qquad
 |f_2(e_2,t)-K_0|\le e_2,
\]
where $P_\xi=\sum_i v_i(B_i(0))_+\in[P_-,P_+]$.
Consequently $|f_2(e_2,t)|\le K_++e_2<0.011$ on this entire interval.
\end{lemma}
\begin{proof}
The comparison rows are $c^\dagger=Av+a(0)_-$ and $D^\dagger=B(0)$,
with input row $B(0)$. Their output at $e_2$ is exactly
$(P_\xi(A-A_0)+C_0,K_0)$, including initially negative neurons.
The error in the first output coordinate is at most
$\|c-c^\dagger\|\|(B(0))_+\|+\|c\|\|B-B(0)\|$.
The second coordinate obeys the corresponding bound with
$D-B(0)$ and $\|D\|$. Lemma~\ref{sbsharp:weighted}, $S<2$ and
\eqref{eq:good-moments} give the displayed $e_1,e_2$.
No probe gates are held fixed.
\end{proof}


\subsection{Uniform initial OOD descent}
\begin{lemma}[Instantaneous residual-kernel lower bound]\label{lem:kernel-lower}
Let $\xi,v$ be unit vectors, $c=\xi^\top v\ge0$, and
$L=\sum_i(u_i^\top\xi)_+(u_i^\top v)_+$.
Suppose $\|r_\xi-\gamma\xi\|,\|r_v-\gamma v\|\le d$, with $\gamma\ge0$.
Under balance, the kernel of Lemma~\ref{lem:probe-kernel} satisfies
\begin{align}\label{eq:kernel-lower}
 \langle r_\xi,\mathsf K_t(\xi,v)r_v\rangle
 \ge{}&2\gamma^2cL-(L+cS)(2\gamma d+d^2)\nonumber\\
 &-c(\|W\|_F+\|U\|_F)\|W-U\|_F(\gamma+d)^2.
\end{align}
\end{lemma}
\begin{proof}
The scalar-identity part of the kernel has leading pairing
$\gamma^2cL$ and residual cost at most $L(2\gamma d+d^2)$.
In the second part replace $w_iw_i^\top$ by $u_iu_i^\top$ while retaining
its actual selectors. Its operator-norm change is bounded by
$c(\|W\|_F+\|U\|_F)\|W-U\|_F$.
The leading pairing is again $\gamma^2cL$, since
$c_i(\xi)\alpha_i(v)(u_i^\top\xi)(u_i^\top v)
=(u_i^\top\xi)_+(u_i^\top v)_+$, including zero preactivations.
The second residual cost is at most $cS(2\gamma d+d^2)$.
Combining the terms proves \eqref{eq:kernel-lower}.
\end{proof}

Put $\gamma=1-s/4=0.99$ and define locally
\begin{equation}\label{eq:lambda-local}
 \lambda_0=(\gamma+d_{\rm train})
 \frac{(1+a_D)^2+(\mu_2+a_D)^2}{2}.
\end{equation}
Indeed $2\mathcal L(0)\le(\gamma+d_{\rm train})^2\mathbb E_n\|x\|^2$,
so Lemma~\ref{sbwindow:mass} applies with this rate. At $h=10^7$,
\[
 \lambda_0=0.49501865088\ldots<0.495019;
\]
it decreases as $h$ increases. Use $F,D_*$ from Lemma~\ref{sbwindow:mass}.
For a unit-probe arc with initial radial error at most $d_0$, strong-cluster
inner product at least $c_0>0$, and initial ReLU-product sum at least $L_0$,
set
\[
 d=d_0+F,\quad L_-=L_0-F,\quad S_+=se^{2\lambda_0t},
\]
\begin{equation}\label{eq:descent-budget}
 \begin{split}
 \mathcal J(t)={}&2\gamma^2c_0L_-
 -(L_-+c_0S_+)(2\gamma d+d^2)
 -c_0D_*(\gamma+d)^2.
 \end{split}
\end{equation}


\begin{lemma}[A finite initial-descent interval]
\label{sbsharp:descent}
Suppose the right-hand side of \eqref{eq:kernel-lower}, with
$\xi$ and each normalized weak sample as its two directions, has a
nonnegative lower bound obtained from the same envelopes.
While $L_->0$, $2\gamma^2c_0>2\gamma d+d^2$, and
$2\gamma^2L_--S_+(2\gamma d+d^2)-D_*(\gamma+d)^2>0$,
\[
 \dot E(\xi_\theta,t)\le-\frac{(1-a_D)^2}{2}\mathcal J(t).
\]
The function $\mathcal J$ is strictly decreasing until its first zero
$T_0$ in this regime. For every $0<T<T_0$ the uniform descent rate and
integrated drop are
\[
 c(T)=\frac{(1-a_D)^2}{2}\mathcal J(T)>0,\qquad
 a(T)=\frac{(1-a_D)^2}{2}\int_0^T\mathcal J(t)\,dt>0.
\]
Thus $T_0$ is the supremum of the strictly negative-derivative intervals
certified by \emph{this budget}, not a claimed optimal interval for the
actual flow. The integral is an explicit polynomial antiderivative in
$x=e^{\lambda_0t}$: integrate its constant term as $T$ and each $x^j$
term as $(e^{j\lambda_0T}-1)/(j\lambda_0)$.
\end{lemma}
\begin{proof}
Apply Lemma~\ref{lem:kernel-lower} and the displacements in
Lemma~\ref{sbwindow:mass}. Positive homogeneity contributes the squared
sample norm for each normalized training direction. The two displayed positivity conditions
make that bound increasing in both $c$ and $L$, permitting their lower
bounds to be substituted without replacing a negative term incorrectly.
Multiply by the strong-cluster weight and its squared sample norm.
For monotonicity, write the first two terms as
$L_-(2\gamma^2c_0-2\gamma d-d^2)-c_0S_+(2\gamma d+d^2)$.
The first product decreases while its factors are positive, and both
subtracted terms increase. Integration proves the claims.
\end{proof}

\subsubsection{The deterministic probe arcs}


Use Lemma~\ref{sbsharp:descent} with the local rate \eqref{eq:lambda-local},
$c_0=\sin(3.5\cdot10^{-5})-2a_D$,
$L_0=s(1/(4\pi)-0.002)\cos(4.5\cdot10^{-5})-2a_Ds$,
and
\[
 d_0=\sqrt2s\left\{\sqrt{\frac{\log10400}{4h}}
        +\frac{4\log10400}{3h}\right\}+0.05(5\cdot10^{-7}).
\]
Its worst value is less than $2.729707\cdot10^{-5}$.
For this inner-arc application set $\beta=4.5\cdot10^{-5}$.
Use $c_w=\cos(4.5\cdot10^{-5})-2a_D/\mu_2$,
$L_{w,0}=0.248s-s(4.5\cdot10^{-5}+2a_D/\mu_2)$,
and $d_{w,0}=\max\{d_0,d_{\rm train}\}$ in
\eqref{eq:descent-budget}, replacing $c_0,L_0,d_0$ by these weak
bounds. This budget stays positive.
At $h=10^7$, the strong budget's zero belongs to
$(0.00018270,0.00018271)$, and at $T=0.00017$ its weighted rate exceeds
$1.5604\cdot10^{-9}$. The exact polynomial antiderivative gives a
drop greater than $2.04297\cdot10^{-12}$, proving the reported bounds.
All substitutions respect the two monotonicity conditions in the lemma.
For larger $h$, $d_{\rm net}$ and $\lambda_0$ decrease; the resulting
residual and displacement charges decrease. The rate and integrated-drop
lower bounds therefore hold uniformly for $h\ge10^7$.
To extend to the larger arc, anchor the initial residual at
$\alpha=3.5\cdot10^{-5}$. Its Lipschitz constant $0.05$ gives the
pointwise radial budget $d_0+0.05(\theta-\alpha)$, while use
$c(\theta)=\sin\theta-2a_D$ and
$L_0(\theta)=s(1/(4\pi)-0.002)\cos\theta-2a_Ds$.
For $\alpha\le\theta\le4\cdot10^{-4}$ and $t\le0.00017$, direct
partial differentiation gives
$\mathcal J_c>0.0060$, $|\mathcal J_d|<0.0063$, and
$|\mathcal J_L L_0'(\theta)|<2\cdot10^{-9}$.
Thus $\partial_\theta\mathcal J>0.0060\cos(4\cdot10^{-4})
-0.05(0.0063)-2\cdot10^{-9}>0.005$.
The weakest strong budget remains the left endpoint's, so the same
drop and rate apply to the larger arc without any new probability cost.
For the full arc set $\beta=4\cdot10^{-4}$ explicitly:
$c_w=\cos(4\cdot10^{-4})-2a_D/\mu_2$,
$L_{w,0}=0.248s-s(4\cdot10^{-4}+2a_D/\mu_2)$, and
$d_{w,0}=\max\{d_{\rm train},d_0+0.05(4\cdot10^{-4}-\alpha)\}$.
The weak budget remains strictly positive through $0.00017$.

For the longer initial transient on the upper subarc, keep the same
probability event and use the \emph{uniform} radial budget
$d_0+0.05(4\cdot10^{-4}-3.5\cdot10^{-5})$, cosine at
$4\cdot10^{-4}$ in $L_0$, and $c_0=\sin(3\cdot10^{-4})-2a_D$.
This is merely Lipschitz transport from the existing anchor.
At $h=10^7$, the corresponding budget zero lies in $(0.00571064,0.00571066)$.
At $T=0.0055$, direct polynomial integration gives a drop greater than
$2.15495\cdot10^{-9}$ and a uniform derivative magnitude greater than
$2.6743\cdot10^{-8}$. For this upper subarc use exactly the full-arc weak bounds with
$\beta=4\cdot10^{-4}$ just defined; their budget is positive through
$0.0055$. Thus the longer
transient and the later rebound coexist on this same $10^{-4}$-wide arc.



\subsection{Uniform post-transient increase}\label{sec:late-increase}

For every time with
$A^2\in[0.69,0.95]$, the following calculation applies.
Put $\bar\theta=4\cdot10^{-4}$, $A_-^\dagger=\sqrt{0.69}$,
$k_-=0.05-M_*$, and
\begin{align*}
 T_-&=A_-^\dagger P_- -3q_*\sqrt{K_+}-PH,\\
 T_+&=\sqrt{0.95}P_++3q_*\sqrt{K_+}+PH,\\
 F_-&=P_-(A_-^\dagger-A_{0,+})-C_+-e_1-3\bar\theta,\\
 V_1&=k_-\cos\bar\theta\,T_-
   -\bar\theta\{l_*+\sqrt2(l_*+\beta_BZ_*+P)\}-2\sqrt2P,\\
 V_2&=Z_*(T_++\bar\theta)+\bar\theta(T_++D_*^{\rm mov})
   +\bar\theta\beta_B(l_*+\beta_BZ_*+P)+P(\sqrt2+\beta_B).
\end{align*}
Using Lemma~\ref{lem:gf-regularity}, for
$q_i=a_i\sin\theta+B_i\cos\theta$, differentiate the actual
all-neuron sum $(c,D)(q_i)_+$. The leading first-coordinate term is
$k\sum a_i^+(q_i)_+\ge k\cos\theta\sum a_i^+(B_i)_+$.
The other potentially negative longitudinal leading term satisfies
$k\sin\theta\sum c_i a_i^+\mathbf1_{q_i>0}\ge-\sin\theta\,l_*$,
because $c_i=a_i+p_i$. The remaining input terms cost at most
$\sin\theta\sqrt2(l_*+\beta_BZ_*+P)$, and the two output/input error
terms cost $2\sqrt2P$. Thus $f_1'\ge V_1>0$ at the parameters certified below.
For the second coordinate, retain the leading signed term
$-R\sum a_i^+(q_i)_+$ and then bound its magnitude by
$Z_*(T_++\bar\theta)$. For the input-leading term, preserve its sign:
\[
 \sum D_i a_i^+\mathbf1_{q_i>0}
 \le\sum (B_i(0))_+a_i^++\|D-B(0)\|u
 \le T_++D_*^{\rm mov}.
\]
Its upper bound is consequently $\bar\theta(T_++D_*^{\rm mov})$.
The remaining terms cost
$\bar\theta\beta_B(l_*+\beta_BZ_*+P)$, and the two error terms cost
$P(\sqrt2+\beta_B)$. Hence $f_2'\le V_2$; an absolute bound is neither
claimed nor needed. Since $f_2-\cos\theta<0$, this one-sided bound has
the correct sign in the error derivative.
These formulas permit both training and probe gate changes.
The output envelope and spatial Lipschitz bound
$|\partial_\theta(f_1-\sin\theta)|\le S+1<3$ give
$f_1-\sin\theta\ge F_->0$.
For the second coordinate,
\[
 0.248s-e_2-\beta_B\sqrt2\,\bar\theta
 \le f_2(\xi_\theta,t)
 \le K_++e_2+\beta_B\sqrt2\,\bar\theta.
\]
The endpoints lie in $(0,0.012)$ and $\cos\theta>0.99$, so
$f_2-\cos\theta<0$ and $|f_2-\cos\theta|<1.011$. Therefore, in physical time,
\[
 \dot E\ge\tfrac12(F_-V_1-1.011V_2)>3.5618\cdot10^{-6}.
\]
All constants here are the exact definitions, bounded by the numerical
ledger in the proof of Theorem~\ref{sbsharp:theorem}. The physical endpoint
separation is
$\log\frac{(0.95)(0.285)}{(0.05)(0.715)}>2.0246$; integration gives a
full-arc rebound exceeding $7.2\cdot10^{-6}$. Substitution of
$\bar\theta=4.5\cdot10^{-5}$ gives a rate above $10^{-5}$ and an
inner-arc rebound above $2\cdot10^{-5}$. Absolute continuity gives strict increase on $[t_-,t_b]$.


\subsection{Assembly and minimum location}

\begin{proof}[Proof of Theorem~\ref{sbsharp:theorem}]

Proposition~\ref{prop:good-event} supplies one common event of probability
at least $0.97$. Formula~\eqref{eq:logistic-time} establishes the time
ordering and $t_b<7.6$.

Set $H=3.8$, $\nu=0.05$, $A_{0,\pm}=\sqrt{0.04(0.25\pm0.002)}$.
The normalized terminal time is at most $3.773836<H$. Direct
substitution into Lemma~\ref{sbsharp:weighted} gives the following
outward-rounded upper bounds (the minimum-width case is worst):
\[
\begin{array}{c|c@{\quad}c|c}
 b_0&7.906\cdot10^{-5}&P&6.001\cdot10^{-7}\\
 \beta_B&0.200003&Z_*&1.316\cdot10^{-4}\\
 l_*&1.021\cdot10^{-4}&W_*&2.337\cdot10^{-4}\\
 q_*&5.337\cdot10^{-4}&C_c^*&0.002094\\
 D_*^{\rm mov}&1.990\cdot10^{-4}&M_*&0.002602\\
 V_*&0.010777&J_D&0.009034\\
 J_B&2.315\cdot10^{-5}&J_p&9.422\cdot10^{-4}\\
 \Gamma_C&0.016426&&
\end{array}
\]
All stopping conditions are strict. The mass-stop expression is below
$1.056$; the learned core mass exceeds $0.55196$, its angle is below
$0.12210$, and its normalized share exceeds $0.80387$.
The calculation uses the exact definitions, not successive substitution
of the displayed rounded values.



The share used here is
$\omega_C=(1/8+1/(4\pi)-0.002)/0.252$.
Lemma~\ref{sbwide:core-holding}, with $A_{\rm spec}=\sqrt{0.71}$, proves the mass and angular statements
on $[t_{\rm spec},t_b]$. Lemmas~\ref{sbsharp:weighted} and
\ref{lem:output-envelope} give $S<2$ and the weak-fitting bound on
$[0,t_b]$. Lemma~\ref{sbsharp:descent} with the three explicit arc
budgets proves initial descent. Section~\ref{sec:late-increase} proves
both rebound bounds on those same arcs.
For each probe, absolute continuity gives continuity on $[0,t_b]$, so a
global minimum exists. Strict initial descent excludes zero. Strict
increase on $[t_-,t_b]$ excludes every time in $(t_-,t_b]$.
Thus every global minimum lies in $(0,t_-]$, strictly before
$t_{\rm spec}$. The sector has width $4\cdot10^{-4}-3.5\cdot10^{-5}$.
Every probe is more than $0.9$ from every training sample on the same
sample event. All conclusions therefore coexist on the single event.
\end{proof}

\section{Numerical validation of the certified theorem regime}\label{app:validation}

\subsection{Setup}
The exact smallest-width
point of Theorem~\ref{sbsharp:theorem} was instantiated with $h=10^7$,
$\varepsilon=\sqrt{0.04/h}$, $\mu_1=1$, $\mu_2=10^{-4}$,
$\sigma_1=\sigma_2=10^{-9}$, all $2000$ noisy training samples,
and random seed $20260918$. Its initialization and sample events passed
without resampling. The empirical gradients were evaluated using the
actual masks: neurons with a mask uniform across a cluster were summed
as matrices, and all boundary neurons were evaluated sample by sample.
An all-pairs preflight containing both types of boundary neuron passed
the driver's $10^{-13}$ output-and-gradient agreement tolerance.

\subsection{Results}
Classical RK4 runs with maximum physical steps $0.02$ and $0.01$,
smaller initial steps, and exact snapshot stops gave the following
\emph{numerical} values:
\[
\begin{array}{c|r}
 t_{\rm spec}&5.489569515\\
 t_a&5.513978830\\
 t_b&7.538624447\\
 \text{core mass at }t_{\rm spec}&0.581020472\\
 \text{largest core angle there}&0.118179854\\
 \max S&0.979990746\\
 \max |f_2(e_2)|&0.009992689
\end{array}
\]
For $\theta=4\cdot10^{-5}$, the observed secant was
$1.05870369\cdot10^{-4}$. A sampled derivative crossing occurred between
$0.04617$ and $0.04717$, and the sampled maximum initial improvement was
$2.76713\cdot10^{-9}$.
For $\theta=3.5\cdot10^{-4}$, the observed secant was
$1.04807195\cdot10^{-4}$, the sampled crossing bracket was
$[0.38017,0.40017]$, and the sampled improvement was
$2.36472\cdot10^{-7}$. Both are early transients well before the
specialization declaration.

\subsection{Numerical consistency and reproducibility}
The inner-probe trajectory was integrated with maximum RK4 steps $0.02$
and $0.01$. Both runs use the same imposed early-time caps:
$10^{-5}$ through $t=0.00017$ and $10^{-3}$ through $t=0.1$.
At the $492$ snapshots common to the two runs, halving the maximum step
changed the inner-probe error by at most $5.04\cdot10^{-14}$, the core
learned mass by at most $9.10\cdot10^{-11}$, the total mass by at most
$1.12\cdot10^{-10}$, and the error derivative by at most
$4.01\cdot10^{-14}$. The two runs share the same early-window step
sequence, so this refinement comparison is informative only after
$t=0.1$; in particular, it is not an integration-error estimate for
either certified initial-descent interval. These are numerical consistency
checks, not rigorous integration-error bounds, continuous-time certificates,
minimum-location proofs, or proof premises.

The analysis script evaluates $22$ theorem-facing numerical conditions
covering mass, angles, relative residual, mismatch, comparison errors,
core holding, and initial/late derivatives and aborts unless all are true.
All $22$ assertions pass for the archived fine/upper trajectories; the
coarse trajectory supplies the refinement comparison. The coarse and
upper files are bit-identical on all recorded training-state columns, so
their two probe outputs are observations of one numerical training
trajectory. The fixed-seed driver uses no rejection sampling: if the
prescribed initialization/sample event fails, execution terminates rather
than retrying another seed.

The accompanying validation archive contains the driver, analysis script,
raw trajectories, initialization manifests, arithmetic certificate ledger,
summary, and provenance record. The coarse trajectory was regenerated
directly from the archived unchanged C++ source with maximum step $0.02$;
its initialization manifest is unchanged, all source step caps are
satisfied, and its timestamp grid is a subset of the fine grid.
The summary and provenance records contain SHA-256 digests for the archived
source and output artifacts. These hashes certify the integrity of the
archived files; they are not a claim that recompiling under a different
compiler, optimization pipeline, or hardware architecture reproduces
floating-point outputs byte-for-byte. The provenance record therefore also
stores the compiler, architecture, compile command, run command, source
hash, and binary hash. Figure~\ref{fig:certified-validation-app} summarizes
the archived numerical validation; it is not used as a proof premise.

\begin{figure}[t]
\centering
\includegraphics[width=\linewidth]{SWINGBY_STRENGTHENING_VALIDATION.png}
\caption{Numerical validation at the smallest-width certified parameter
point. Top: OOD error changes and derivatives for the inner and upper
probes, including their early transients and later increases. Bottom:
coherent-core mass and angle together with the weak-concept output. The
dashed line marks the declared specialization time and the shaded interval
marks the theorem's rebound snapshots. The figure is numerical validation
only and is not a proof premise.}
\label{fig:certified-validation-app}
\end{figure}


\section{Exact noiseless examples}\label{app:noiseless}
The leakage scale in this appendix is inherited from initialization.
The results concern one or two neurons and are separate from
Theorem~\ref{sbsharp:theorem}.

Consider the planar ReLU autoencoder
\[
 f(x)=\sum_i w_i(u_i^\top x)_+,\qquad
 \mathcal L=\frac14\bigl(\|f(\mu_1e_1)-\mu_1e_1\|^2
                         +\|f(\mu_2e_2)-\mu_2e_2\|^2\bigr),
 \quad \mu_1,\mu_2>0.
\]
Time is gradient-flow time for this normalization. Probe error is
$E(\xi,t)=\|f_t(\xi)-\xi\|^2/2$.

\begin{theorem}[Uniform sector reversal]
\label{sbsector:theorem}
Take one neuron and
\[
 u(0)=w(0)=(a_0,-\eta),\qquad
 a_0,\eta>0,\qquad a_0^2+\eta^2\le\frac1{16}.
\]
Write $\lambda=\mu_1^2/2$, $\beta=1-\eta^2$, and, for every
$z\in[1/2,1]$, let
\[
 k_z=\sqrt{1-\eta^2z^2},\qquad \xi_z=(\eta z,-k_z).
\]
The following conclusions hold for the actual gradient flow.
\begin{enumerate}[label=(\roman*)]
\item Both training gates and all these probe gates are fixed for every
finite time. There are increasing functions $a,c$ and a decreasing
positive function $D$ such that
\[
 u=(a,-\eta),\quad w=(c,-\eta D),\qquad
 M=wu^\top=
 \begin{pmatrix}ac&-\eta c\\-\eta aD&\eta^2D\end{pmatrix}.
\]
The same matrix $M$ acts on the entire probe sector. Its four scaled
coefficients are bounded by $2$, their time derivatives by $5\lambda$,
and transverse input/output energy is at most $2\eta^2$.

\item Define the finite hitting times
\[
 t_-:=\inf\{t:a(t)=1/4\},\qquad
 t_+:=\inf\{t:a(t)=1/\sqrt2\}.
\]
For every $z\in[1/2,1]$, $\dot E(\xi_z,t)<0$ on $[0,t_-]$ and
$\dot E(\xi_z,t)>0$ on $[t_+,\infty)$. Every global minimizing time
$t_{\rm rev}(z)$ is a strict local minimum and satisfies
\begin{equation}\label{sbsector:time}
 \frac1\lambda\log\frac1{4a_0}
 \le t_-<t_{\rm rev}(z)<t_+
 \le\frac1{2\lambda\beta}
 \log\frac{\beta-a_0^2}{a_0^2(2\beta-1)}.
\end{equation}
No uniqueness of the minimum is asserted.

\item If $E_\infty(\xi_z)=\lim_{t\to\infty}E(\xi_z,t)$, then
\begin{equation}\label{sbsector:amplitude}
 \frac14\eta^2
 < E_\infty(\xi_z)-\min_{t\ge0}E(\xi_z,t)
 <4\eta^2.
\end{equation}
The sector has angular width
\[
 \arcsin\eta-\arcsin(\eta/2),\qquad
 \frac\eta2\le\arcsin\eta-\arcsin(\eta/2)\le\eta.
\]
Thus its width is $\Theta(\eta)$ and its rebound amplitude is
$\Theta(\eta^2)$, with constants uniform in $a_0,\eta,z$ in the stated
range. The upper and lower times in \eqref{sbsector:time} are of order
$\lambda^{-1}\log(1/a_0)$ as $a_0\to0$.
\end{enumerate}
\end{theorem}

\begin{corollary}[Two learned concepts and two reversal sectors]
\label{sbsector:two-concept}
Take two neurons initialized by
\[
 u_1(0)=w_1(0)=(a_0,-\eta),\qquad
 u_2(0)=w_2(0)=(-\zeta,b_0),
\]
where $a_0,b_0,\eta,\zeta>0$ and
$a_0^2+\eta^2\le1/16$, $b_0^2+\zeta^2\le1/16$.
Both training reconstructions converge to their targets. The conclusions
of Theorem~\ref{sbsector:theorem} hold for the strong probes
$(\eta z,-\sqrt{1-\eta^2z^2})$, $z\in[1/2,1]$.
They also hold, with $(a_0,\eta,\lambda)$ replaced by
$(b_0,\zeta,\mu_2^2/2)$, for the weak probes
$(-\sqrt{1-\zeta^2z^2},\zeta z)$.
The sector widths and rebound amplitudes are respectively
$\Theta(\eta),\Theta(\eta^2)$ and
$\Theta(\zeta),\Theta(\zeta^2)$.
\end{corollary}

\begin{corollary}[Isotropic initialization with positive probability]
\label{sbsector:isotropic}
In the same noiseless planar dataset, take $h=2$ and independently
initialize $u_i(0)=w_i(0)=\varepsilon\widehat u_i$, with
$\widehat u_i$ uniform on $\mathbb S^1$ and $0<\varepsilon\le1/4$.
With probability $1/8$, one neuron starts in quadrant IV and the other
in quadrant II. On that event, both concepts are fitted and
Corollary~\ref{sbsector:two-concept} holds with the positive coordinate
magnitudes supplied by initialization, after relabeling if necessary.

There is a subevent of probability $1/72$ on which both learned
concepts have a probe sector of width between $\varepsilon/4$ and
$\varepsilon$, and every probe in each sector has rebound amplitude
strictly between $\varepsilon^2/16$ and $3\varepsilon^2$.
For each concept $p$ and every such probe, every global minimizing time
satisfies the deterministic bounds
\begin{equation}\label{sbsector:isotropic-time}
 \frac1{\lambda_p}\log\frac1{2\sqrt3\varepsilon}
 <t_{\rm rev}
 <\frac8{15\lambda_p}\log\frac{32}{7\varepsilon^2},
 \qquad \lambda_p=\mu_p^2/2.
\end{equation}
The probability constants are lower guarantees from explicit angular
events; no claim is made about the remaining initializations.
\end{corollary}

\subsection{Exact gradient flow and invariant region}
As long as the first training gate is on and the second is off,
the exact equations, starting with $a(0)=c(0)=a_0$, $D(0)=1$, are
\begin{equation}\label{sbsector:ode}
 \dot a=\lambda[c-a(a^2+\eta^2)],\qquad
 \dot c=\lambda a(1-ac),\qquad
 \dot D=-\lambda a^2D.
\end{equation}
For example, the strong-sample output is $\mu_1a(c,-\eta D)$;
the input-gradient equation and the balance identity below give the
first equation. The second input coordinate is the constant $-\eta$.
In the reduced equations the derivative of
$c^2+\eta^2D^2-a^2-\eta^2$ is $-2\lambda a^2$ times that quantity.
It vanishes initially, giving balance:
\begin{equation}\label{sbsector:balance}
 c^2+\eta^2D^2=a^2+\eta^2.
\end{equation}
Let $a_\infty>0$ be defined by
\[
 a_\infty^2(a_\infty^2+\eta^2)=1.
\]
Then $a_0<a_\infty<1$ and $a_\infty>1/\sqrt2$.
Set $h=c-a(a^2+\eta^2)$. Initially $h=a_0(1-a_0^2-\eta^2)>0$.
On $h=0$ with $0<a<a_\infty$,
\[
 \dot h=\lambda a[1-a^2(a^2+\eta^2)]>0.
\]
Thus $h>0$ while $a<a_\infty$. At any hypothetical first finite hit
of $a_\infty$, $D>0$ and balance would give
$c<\sqrt{a_\infty^2+\eta^2}=1/a_\infty$, hence $h<0$, a contradiction.
Consequently
\[
 0<a_0\le a<a_\infty<1,\quad
 a\le c\le\sqrt{a^2+\eta^2},\quad ac<1,\quad
 \dot a>0,\quad\dot c>0,\quad 0<D\le1.
\]
These bounded quantities extend globally; $a>0$ keeps the first gate on,
and the constant negative second input coordinate keeps the second off.
Since $a\ge a_0$, $D\to0$. Monotonicity, balance, and
\eqref{sbsector:ode} imply $a\to a_\infty$ and $c\to1/a_\infty$.

For every probe under consideration,
$u^\top\xi_z=\eta(az+k_z)>0$. Thus its exact matrix is the matrix in
the theorem. The bounds $a\le1$, $c\le2$, $D\le1$,
$\dot a\le2\lambda$, $\dot c\le\lambda$, and
$|\dot D|\le\lambda$ prove the coefficient and rate bounds by the
product rule. Transverse energy equals $\eta^2(1+D^2)$.

\subsection{Exact error identity and sign estimates}
For this proof abbreviate $k=k_z$ and set
\[
 \mathcal A=c(az+k)-z,
 \qquad \mathcal T=-D(az+k),
 \qquad B=k+\eta^2\mathcal T.
\]
There is no ambient output component. Exactly,
\begin{equation}\label{sbsector:error}
 E=\frac{\eta^2}{2}\mathcal A^2+\frac12B^2,
 \qquad
 \frac{\dot E}{\eta^2}
 =\mathcal A\dot{\mathcal A}+B\dot{\mathcal T}.
\end{equation}
Throughout the trajectory,
\begin{equation}\label{sbsector:basic}
 k\ge\frac{\sqrt{15}}4>\frac34,\qquad
 B\ge k-\eta^2(1+k)\ge\frac58,
 \qquad\dot{\mathcal A}>0.
\end{equation}
Here $az+k\le2$, $\eta^2\le1/16$, and
$\dot{\mathcal A}=\dot c(az+k)+cz\dot a$.
The leakage derivative is exactly
\begin{equation}\label{sbsector:leakage}
 \dot{\mathcal T}
 =\lambda D\{ka^2+z(2a^3+a\eta^2-c)\}.
\end{equation}

If $a\le1/4$, then
$c\le\sqrt{1/8}<3/8$ and $ac\le3/32$. Hence
\[
 \mathcal A=ck-(1-ac)z\le\frac38-\frac{29}{64}
 =-\frac5{64}.
\]
Using $c\ge a$, $z\ge1/2$, and $\eta^2\le1/16$ in
\eqref{sbsector:leakage} also gives
\[
 \dot{\mathcal T}
 \le\lambda Da\{a-z(1-2a^2-\eta^2)\}
 \le-\frac5{32}\lambda Da.
\]
Thus, on $[0,t_-]$,
\begin{equation}\label{sbsector:early-sign}
 \dot E\le-\frac{25}{256}\lambda\eta^2Da<0.
\end{equation}

For $a\ge1/\sqrt2$, the inequality
$c\le a+\eta^2/(2a)$ yields
\[
 2a^3+a\eta^2-c
 \ge(2a^2-1)\left(a+\frac{\eta^2}{2a}\right)\ge0.
\]
Therefore
\[
 \dot{\mathcal T}\ge\lambda Dka^2,
 \qquad
 \mathcal A\ge ak-(1-a^2)
 \ge\sqrt{15/32}-1/2>0.
\]
Equations \eqref{sbsector:error}--\eqref{sbsector:basic} imply, for
every $t\ge t_+$,
\begin{equation}\label{sbsector:late-sign}
 \dot E\ge\frac{15}{32}\lambda\eta^2a^2D>0.
\end{equation}
Unlike an asymptotic comparison, this estimate applies from an explicit
progress level and retains the natural decaying rate $a^2D$.

\subsection{Reached times and a uniform rebound amplitude}
The invariant-region proof guarantees that $t_-$ and $t_+$ are finite.
Since $(a+c)'\le\lambda(a+c)$ and $c\ge a$,
\[
 a(t)\le a_0e^{\lambda t},
 \qquad t_-\ge\lambda^{-1}\log(1/(4a_0)).
\]
In the other direction, $c\ge a$ gives
\[
 \dot a\ge\lambda a(\beta-a^2),\qquad \beta=1-\eta^2.
\]
Comparison with the scalar logistic equation for $a^2$ through $t_+$
gives the upper bound in \eqref{sbsector:time}.

The error is analytic at every finite time because the relevant gates
are fixed and the equations are analytic. It decreases through $t_-$
and increases from $t_+$ onwards. Its global minimum is therefore attained
in $(t_-,t_+)$. A nonconstant real analytic function has isolated local
minima, so every global minimizing time is a strict local minimum.

To bound $D(t_+)$ from below, integrate with respect to the increasing
variable $a$ up to $1/\sqrt2$:
\[
 \frac{d\log D}{da}=-\frac{\lambda a^2}{\dot a}
 \ge-\frac{a}{\beta-a^2},
 \qquad
 D(t_+)\ge
 \sqrt{\frac{\beta-1/2}{\beta-a_0^2}}
 \ge\frac{\sqrt7}{4}.
\]
Integrating \eqref{sbsector:late-sign}, using
$-\dot D=\lambda a^2D$, proves
\[
 E_\infty-E(t_+)
 \ge\frac{15}{32}\eta^2D(t_+)
 \ge\frac{15\sqrt7}{128}\eta^2>\frac14\eta^2.
\]
This also bounds the rebound from the smaller global minimum below.
For the upper bound, expand \eqref{sbsector:error}:
\[
 E-\frac12
 =\eta^2\left(\frac{\mathcal A^2}{2}-\frac{z^2}{2}
                         +k\mathcal T\right)
   +\frac{\eta^4}{2}\mathcal T^2\ge-\frac52\eta^2.
\]
At infinity, $ac\to1$ and $D\to0$, so
\[
 E_\infty=\frac{k^2}{2}(1+\eta^2/a_\infty^2),
 \qquad
 E_\infty-\frac12\le\frac{17}{32}\eta^2,
\]
where $a_\infty^{-2}=a_\infty^2+\eta^2\le17/16$.
Thus the rebound is at most $(97/32)\eta^2<4\eta^2$.
Finally, integrating the derivative of $\arcsin$ on
$[\eta/2,\eta]$, with $\eta\le1/4$, proves the sector-width bounds.
This completes the proof of Theorem~\ref{sbsector:theorem}.

\subsection{Proof of the two-concept extension}
As long as the two neurons keep their pure training gates, the first
sees only $\mu_1e_1$ and the second only $\mu_2e_2$. Their equations
are two independent copies of \eqref{sbsector:ode}, with the coordinates
swapped for the second neuron. The first has input $(a,-\eta)$ and
the second $(-\zeta,b)$ with $a,b>0$ and constant $\eta,\zeta>0$;
the invariant-region argument for \eqref{sbsector:ode} preserves those gates
for all finite time. Each training output converges to its own target.
At a strong probe the second preactivation is
$-\zeta\eta z-bk_z<0$; at a weak probe the first is likewise negative.
Each probe thus has exactly the one-neuron error already analyzed.
This proves Corollary~\ref{sbsector:two-concept} without a neglected
complement contribution.

For Corollary~\ref{sbsector:isotropic}, two prescribed quadrants have
probability $(1/4)^2$ in a prescribed neuron order. The two possible
orders are disjoint, giving $1/8$. For the smaller event restrict the
quadrant-IV angle to $[-\pi/3,-\pi/6]$ and the quadrant-II angle to
$[2\pi/3,5\pi/6]$. Each arc has probability $1/12$, so the two orders
have total probability $1/72$. On these arcs all four relevant coordinate
magnitudes lie in $[\varepsilon/2,\sqrt3\varepsilon/2]$.
Substitution into the width and amplitude bounds gives the claims.
In the upper time bound use
$\beta\ge15/16$, $2\beta-1\ge7/8$, and
$a_0^2\ge\varepsilon^2/4$ (or $b_0^2$ for the other concept).
The lower bound uses $a_0,b_0\le\sqrt3\varepsilon/2$.
This proves \eqref{sbsector:isotropic-time}.


\subsection{Uniform normal-form accuracy along the trajectory}
\begin{proposition}[Uniform normal form and derivative accuracy]
\label{sbsector:profile}
Under the hypotheses of Theorem~\ref{sbsector:theorem}, let
$a,c,D$ be its coefficient trajectories and $\lambda=\mu_1^2/2$. Define
\[
 H_z=\frac12[(ac-1)z+c]^2-D(az+1)-\frac12z^2,
 \qquad G_z=\frac{dH_z}{dt}.
\]
For every $t\ge0$ and $z\in[1/2,1]$,
\[
 \left|E(\xi_z,t)-\frac12-\eta^2H_z(t)\right|\le20\eta^4,
 \qquad
 |\dot E(\xi_z,t)-\eta^2G_z(t)|\le80\lambda\eta^4.
\]
The same bounds hold on each of the two sectors in
Corollary~\ref{sbsector:two-concept}, with its own parameters.
\end{proposition}
\begin{proof}
For the fixed probe $(\eta z,sk)$, where $s\in\{-1,1\}$ and
$k=\sqrt{1-\eta^2z^2}$, take a graded matrix
with entries $(\alpha,\eta b;\eta c_*,\eta^2d)$, put
\[
 A_*=(\alpha-1)z+s b,\qquad T_*=c_*z+s d k.
\]
The exact planar identity is
\[
 E=\frac12+\eta^2\left(\frac12A_*^2-s c_*z-d-\frac12z^2\right)
                  +\eta^4K,
\]
where
\[
 K=dz^2-\frac{s z^2(bA_*-c_*z)}{1+k}
       +\frac12T_*^2+\frac{\eta^2b^2z^4}{2(1+k)^2}.
\]
This follows by squaring the exact residual and using
$k-1=-\eta^2z^2/(1+k)$. In our case
$(\alpha,b,c_*,d)=(ac,-c,-aD,D)$ and $s=-1$, giving $H_z$ above.
Theorem~\ref{sbsector:theorem} bounds each coefficient by $2$ and
its rate by $5\lambda$.
Consequently $|A_*|\le5$, $|\dot A_*|\le10\lambda$,
$|T_*|\le4$, $|\dot T_*|\le10\lambda$, and $1+k\ge7/4$.
Termwise estimation gives
\[
 |K|\le2+\frac{12}{7/4}+8+\frac{1/4}{2(7/4)^2}<20.
\]
Differentiation with the probe fixed yields
\[
 \dot K=\dot d z^2
 -\frac{s z^2(\dot b A_*+b\dot A_*-\dot c_*z)}{1+k}
 +T_*\dot T_*
 +\frac{\eta^2b\dot b z^4}{(1+k)^2},
\]
so
\[
 |\dot K|\le\lambda\left(5+\frac{50}{7/4}+40
                                  +\frac{10/16}{(7/4)^2}\right)<80\lambda.
\]
All scaled coefficient bounds have been proved along the actual
trajectory; they are not additional hypotheses of this proposition.
\end{proof}

For this model, $M=wu^\top$ on the entire probe sector and the reduction
remainder $R$ in Appendix~\ref{app:normal} is identically zero. Substitution
of $(\alpha,b,c_*,d)=(ac,-c,-aD,D)$ and $s=-1$ into
\eqref{eq:normal-H}--\eqref{eq:normal-G} gives $H_z,G_z$.
Thus Proposition~\ref{sbsector:profile} verifies the normal-form accuracy
along an actual initialized trajectory, with explicit constants.
The coefficients may still depend on $a_0,\eta$; a universal
parameter-independent limiting profile is not asserted. The reversal
signs follow from \eqref{sbsector:early-sign} and
\eqref{sbsector:late-sign}, rather than from an approximation at a
vanishing sign margin.

\begin{thebibliography}{9}
\raggedright
\bibitem{davis2020} D.~Davis, D.~Drusvyatskiy, S.~Kakade and J.~D.~Lee.
Stochastic subgradient method converges on tame functions.
\emph{Foundations of Computational Mathematics} 20:119--154, 2020.
\href{https://doi.org/10.1007/s10208-018-09409-5}{doi:10.1007/s10208-018-09409-5}.
\bibitem{shevtsova2011} I.~Shevtsova.
On the absolute constants in the Berry--Esseen type inequalities for
identically distributed summands. 2011.
\href{https://arxiv.org/abs/1111.6554}{arXiv:1111.6554}.
\end{thebibliography}
\end{document}
